
\documentclass[11pt]{article}

\usepackage[margin=1in]{geometry}
\usepackage{amsmath, amssymb, amsthm}
\usepackage{mathtools}
\usepackage{enumitem}
\usepackage{hyperref}
\usepackage{microtype}
\usepackage{booktabs}
\usepackage{multirow}

\hypersetup{
  colorlinks=true,
  linkcolor=blue,
  citecolor=blue,
  urlcolor=blue
}

\title{Behavioral Ontology Correctness:\\
From Syntax Checking to Evidence-Valued Verification\\
via SUMO $\to$ GF $\to$ OSLF $\to$ World Models\\[6pt]
\large (AIrxiv Working Paper, 2026-02-20)}
\author{}
\date{}

\newcommand{\DiamondOp}{\mathbin{\Diamond}}
\newcommand{\BoxOp}{\mathbin{\Box}}
\newcommand{\Red}{\;\Rightarrow\;}
\newcommand{\RedStar}{\;\Rightarrow^{*}\;}
\newcommand{\KIF}{\textsc{kif}}
\newcommand{\SUMO}{\textsc{sumo}}
\newcommand{\GF}{\textsc{gf}}
\newcommand{\OSLF}{\textsc{oslf}}
\newcommand{\WM}{\textsc{wm}}
\newcommand{\PLN}{\textsc{pln}}
\newcommand{\OWL}{\textsc{owl}}
\newcommand{\FOL}{\textsc{fol}}
\newcommand{\HOL}{\textsc{hol}}
\newcommand{\ATP}{\textsc{atp}}
\newcommand{\ev}[2]{\langle #1, #2 \rangle}
\newcommand{\Ev}{\mathrm{Ev}}

\theoremstyle{definition}
\newtheorem{definition}{Definition}[section]
\newtheorem{remark}{Remark}[section]
\newtheorem{example}{Example}[section]
\newtheorem{principle}{Principle}[section]

\theoremstyle{plain}
\newtheorem{theorem}{Theorem}[section]
\newtheorem{proposition}{Proposition}[section]

\begin{document}
\maketitle

\begin{abstract}
Large formal ontologies like the Suggested Upper Merged Ontology (\SUMO{}) contain
tens of thousands of axioms in the Knowledge Interchange Format (\KIF{}).
Despite decades of automated-reasoning work, known bug classes persist: arity mismatches,
type confusion, missing arguments, and---more subtly---axioms whose
\emph{operational consequences} diverge from their intended meaning.
Existing debugging methods (glass-box reasoner diagnostics, white-box test
generation, black-box competency questions, OntoClean metaproperty analysis,
design-pattern templates) each address a layer of correctness but remain
disconnected and leave behavioral adequacy unexamined.

We propose a four-stage pipeline---\SUMO{} $\to$ \SUMO{}-\GF{} $\to$ \GF{}
$\to$ \OSLF{} $\to$ \WM{}---that provides \emph{behavioral} and
\emph{evidential} layers of ontology verification currently absent from the
literature.
Stage~1 maps \KIF{} relations and classes to Grammatical Framework (\GF{})
abstract syntax, where the grammar's type discipline catches arity and sort
errors at parse time.
Stage~2 synthesizes Operational Semantics in Logical Form (\OSLF{}) from the
grammar, automatically deriving modal behavioral types ($\DiamondOp$, $\BoxOp$)
and a Galois connection.
Stage~3 interprets behavioral types as evidence-valued world-model queries,
grounding ontological assertions in PLN Evidence counts $\ev{n^+}{n^-}$ and
flagging ungroundable claims.

The \GF{} $\to$ \OSLF{} $\to$ \WM{} core of this pipeline is implemented in
Lean~4 in the referenced Mettapedia pipeline slice (not the entire repository),
with no \texttt{sorry}/axiom placeholders in that slice.
We further propose a \emph{convergent formalization-driven repair loop} that
uses the pipeline's type, behavioral, and evidential signals to disambiguate
and verify ontology corrections---going beyond detection to principled repair.
We survey the ontology debugging literature, document \SUMO{}'s known issues,
and argue that, in the scoped setting of this draft, the proposed pipeline can
integrate major strengths of prior approaches while adding behavioral-correctness
and repair layers.
\paragraph{Status note.} Implementation-status and count claims in this draft are snapshot-scoped; verify against current repository trackers and build outputs before citing as current state. Lean-proven claims are only those tied to explicit Lean file/theorem references; broader mathematical or historical claims are literature-backed context and not asserted as Lean-proven unless explicitly stated.
\end{abstract}

%% ─────────────────────────────────────────────────────────────
\section{Introduction}
\label{sec:intro}

The Suggested Upper Merged Ontology (\SUMO{}) is one of the largest freely available
formal ontologies, comprising over 20\,000 axioms in the Knowledge Interchange
Format (\KIF{}) across 113 domain files~\cite{pease2002sumo}.
It has been the subject of sustained automated-theorem-proving (\ATP{}) effort,
including annual prizes at the CADE ATP System Competition
(CASC)~\cite{pease-casc-prizes,pease-sutcliffe-2007}.
Despite this, known bugs persist: arity mismatches between relation declarations
and their usage sites, type confusion (e.g., applying \texttt{contraryAttribute}
to a \texttt{Process} where an \texttt{Attribute} is expected), missing arguments
in domain-module facts, and semantically problematic axioms that pass syntactic
checks but produce unintended inferences.

The ontology engineering community has developed several debugging paradigms:
\begin{itemize}[nosep]
  \item \textbf{Glass-box}: Tableau-internal axiom pinpointing for unsatisfiable
    \OWL{} classes~\cite{kalyanpur2005}.
  \item \textbf{White-box}: Automatic test generation from axiom syntax, \ATP{}-based
    redundancy detection~\cite{keet2017whitebox}.
  \item \textbf{Black-box}: WordNet-based competency questions for external
    validation~\cite{keet2017blackbox}.
  \item \textbf{OntoClean}: Philosophical metaproperty analysis (rigidity, identity,
    unity) for taxonomic adequacy~\cite{guarino-welty-ontoclean}.
  \item \textbf{Design patterns}: Modular ontology templates for
    prevention~\cite{gangemi-presutti-patterns}.
\end{itemize}
Each addresses a genuine layer of correctness, but they remain disconnected:
no framework integrates syntactic well-formedness, semantic consistency,
\emph{behavioral adequacy} (do the axioms' operational consequences match
the intended domain dynamics?), and \emph{evidential grounding} (can the
asserted knowledge be supported by observation?).

\paragraph{Thesis.}
We argue that the missing layers are \emph{behavioral} and \emph{evidential}
correctness, and that operational semantics provides the bridge.
Concretely, we propose the pipeline
\[
  \SUMO{} \;\xrightarrow{\text{Stage 1}}\; \SUMO\text{-}\GF{}
  \;\xrightarrow{\text{Stage 2}}\; \OSLF{}
  \;\xrightarrow{\text{Stage 3}}\; \WM{}
\]
where:
\begin{enumerate}[nosep]
  \item \textbf{Stage 1} (\SUMO{} $\to$ \SUMO{}-\GF{}): translates \KIF{}
    into Grammatical Framework abstract syntax, gaining grammar-level
    sort and arity checking.
  \item \textbf{Stage 2} (\GF{} $\to$ \OSLF{}): synthesizes behavioral types
    ($\DiamondOp$, $\BoxOp$, native types) from operational semantics.
  \item \textbf{Stage 3} (\OSLF{} $\to$ \WM{}): interprets behavioral
    properties as evidence-valued world-model queries.
\end{enumerate}

The \GF{} $\to$ \OSLF{} $\to$ \WM{} core is implemented in the referenced
Mettapedia Lean~4 pipeline slice with no \texttt{sorry}/axiom placeholders in
that slice (hundreds of generated grammar rules, proven Galois connection, and
proven evidence-valued semantics in the cited modules).
The \SUMO{} $\to$ \GF{} stage builds on our existing \texttt{sumo-lean4}
\KIF{} checker and HOL emitter, and on Ranta et al.'s typeful ontology
verbalization work~\cite{ranta-typeful-ontologies}.

\paragraph{Contributions.}
\begin{enumerate}[nosep]
  \item A four-level correctness taxonomy for ontologies (\S\ref{sec:correctness}).
  \item A survey of \SUMO{}'s structure, known bugs, and their root causes
    (\S\ref{sec:sumo}).
  \item A systematic review of ontology debugging methods and what each
    misses (\S\ref{sec:debugging}).
  \item The \SUMO{} $\to$ \GF{} $\to$ \OSLF{} $\to$ \WM{} pipeline
    architecture (\S\ref{sec:pipeline}).
  \item A comparison with direct \SUMO{} $\to$ Lean approaches, showing
    the pipeline's advantages (\S\ref{sec:comparison}).
  \item A worked example of a real \SUMO{} bug flowing through the
    pipeline (\S\ref{sec:example}).
  \item A formalization-driven repair methodology with canary theorems
    and behavioral disambiguation (\S\ref{sec:repair}).
\end{enumerate}

%% ─────────────────────────────────────────────────────────────
\section{What Makes an Ontology Correct?}
\label{sec:correctness}

We identify four levels of ontology correctness, each subsuming the
previous:

\begin{definition}[Syntactic well-formedness]
Every axiom conforms to the language grammar: balanced parentheses,
declared relation names, correct arity at each usage site, and
well-sorted arguments matching \texttt{domain}/\texttt{range}
declarations.
\end{definition}

\begin{definition}[Semantic consistency]
The axiom set has at least one model (is satisfiable).
No subset of axioms entails $\bot$.
Equivalently: no class is forced to be empty by the axioms unless
intended (unsatisfiable-class debugging in the sense
of~\cite{kalyanpur2005}).
\end{definition}

\begin{definition}[Taxonomic adequacy]
The subsumption hierarchy respects philosophical metaproperties:
\begin{itemize}[nosep]
  \item \textbf{Rigidity}: a rigid property cannot cease to hold for
    an individual (e.g., \textit{Human} is rigid; \textit{Student}
    is anti-rigid).
  \item \textbf{Identity}: a sortal carries an identity criterion
    that its subsorts inherit.
  \item \textbf{Unity}: instances of a unitary property share the
    same ``whole-part'' structure.
\end{itemize}
This is the domain of OntoClean~\cite{guarino-welty-ontoclean}.
A taxonomy is \emph{adequate} if no anti-rigid class subsumes a rigid
one, no class carries conflicting identity criteria, etc.
\end{definition}

\begin{definition}[Behavioral correctness]
\label{def:behavioral}
The operational consequences of the axioms match the intended domain
dynamics.
Formally: given an operational semantics $\Red$ derived from the
ontology's axioms, for every intended domain scenario $\sigma$:
\begin{enumerate}[nosep]
  \item If the domain intends ``$\phi$ is possible from state $s$,''
    then $s \models \DiamondOp\,\phi$ (the operational semantics
    produces a reduction path to a state satisfying~$\phi$).
  \item If the domain intends ``$\phi$ is necessary for state $s$,''
    then $s \models \BoxOp\,\phi$ (every predecessor of~$s$ satisfies
    $\phi$).
  \item No axiom produces zero evidence under world-model evaluation
    unless it is explicitly a vacuous assertion.
\end{enumerate}
\end{definition}

The first two levels are well-served by existing tools (\ATP{} for
consistency, \KIF{} checkers for syntax).
Level~3 (OntoClean) is mature but manual and requires philosophical
expertise.
In the source set surveyed for this draft, we did not find an existing
ontology-engineering methodology that addresses Level~4 behavioral correctness
as an end-to-end verification layer.
Table~\ref{tab:levels} summarizes the landscape.

\begin{table}[ht]
\centering
\caption{Ontology correctness levels and which methods address them.}
\label{tab:levels}
\small
\begin{tabular}{@{}lccccccc@{}}
\toprule
 & \rotatebox{70}{KIF checker} & \rotatebox{70}{ATP (E/Vampire)} & \rotatebox{70}{Glass-box} & \rotatebox{70}{OntoClean} & \rotatebox{70}{Keet testing} & \rotatebox{70}{Axiom weak.} & \rotatebox{70}{\textbf{Pipeline}} \\
\midrule
Syntactic well-formedness & \checkmark & & & & \checkmark & & \checkmark \\
Semantic consistency & & \checkmark & \checkmark & & (\checkmark) & \checkmark & \checkmark \\
Taxonomic adequacy & & & & \checkmark & & & (\checkmark) \\
Behavioral correctness & & & & & & & \checkmark \\
\midrule
Repair guidance & & & (\checkmark) & & & \checkmark$^*$ & \checkmark \\
\bottomrule
\multicolumn{8}{@{}l}{\footnotesize $^*$\,$\mathcal{EL}$ only (Troquard et al.~2018, Baader \& Kriegel 2023).}
\end{tabular}
\end{table}

\begin{remark}
Level~4 (behavioral correctness) presupposes an operational reading of the
ontology.
This is natural for ontologies that describe processes, events,
causation, and change---a large fraction of any upper ontology including
\SUMO{}.
For purely classificatory ontologies (taxonomies without process
axioms), levels~1--3 may suffice.
\end{remark}

%% ─────────────────────────────────────────────────────────────
\section{SUMO: Structure and Known Issues}
\label{sec:sumo}

\subsection{KIF Syntax and Semantics}

The Knowledge Interchange Format~\cite{hayes-menzel-kif-2001} uses a
LISP-compatible S-expression syntax.
Key features distinguishing \KIF{} from standard first-order logic:
\begin{itemize}[nosep]
  \item \textbf{Free syntax}: no syntactic distinction between object
    and predicate variables.
    A variable \texttt{?R} can appear in predicate position:
    \texttt{(?R ?X ?Y)}.
  \item \textbf{Variadic relations}: a relation need not have a unique
    arity.
    \texttt{@ROW} variables allow variable-length argument lists.
  \item \textbf{Row variables}: \texttt{@ROW} splices a sequence of
    terms into an argument list, taking the language beyond first-order.
  \item \textbf{Self-application}: a relation can appear as an argument
    to itself or to \texttt{instance}/\texttt{subclass}.
\end{itemize}

Hayes and Menzel~\cite{hayes-menzel-kif-2001} gave a formal semantics
for a cleaned-up variant called \textsc{skif}, noting that the full use
of row variables takes the language beyond first-order expressiveness.

\subsection{SUMO Architecture}

\SUMO{} is organized into:
\begin{itemize}[nosep]
  \item \textbf{Merge.kif} (607\,KB): the upper ontology---classes like
    \texttt{Entity}, \texttt{Physical}, \texttt{Process},
    \texttt{Attribute}, and foundational relations
    (\texttt{instance}, \texttt{subclass}, \texttt{domain},
    \texttt{range}).
  \item \textbf{Mid-level-ontology.kif} (1.1\,MB): intermediate
    concepts (countries, organizations, artifacts, biological
    processes).
  \item \textbf{Domain modules} (111 files): specialized knowledge
    (Cars, Geography, Economy, Languages, etc.).
  \item \textbf{domainEnglishFormat.kif} (3.7\,MB): English natural
    language generation templates for every relation.
\end{itemize}

Total: 113 \KIF{} files, approximately 1.6\,GB including
natural-language annotations.

\subsection{Known Bug Classes}

We have encountered or documented the following bug classes through
our \texttt{sumo-lean4} checker and manual
inspection\footnote{See also \texttt{PossibleBugsInSUMO\_HopefullySolvedByNowBySomeone.kif}
in our repository.}:

\paragraph{1. Arity mismatches.}
The predicates \texttt{classmate} and \texttt{colleague} are
explicitly ternary (person, person, context), but were also declared as
instances of \texttt{IrreflexiveRelation}---a class defined in
\texttt{Merge.kif} to apply only to binary relations.
Our checker flagged this; the fix was to remove the
\texttt{IrreflexiveRelation} instances and add explicit ternary
irreflexivity axioms:
\begin{verbatim}
  (forall (?A ?COURSE)
    (not (classmate ?A ?A ?COURSE)))
\end{verbatim}

\paragraph{2. Missing arguments.}
In \texttt{mondial.kif}, four Factbook-style demographics facts
were missing the region argument:
\texttt{beliefGroupPercentInRegion} and
\texttt{ethnicityPercentInRegion} had two arguments instead of three.
The checker normalized them to the correct three-argument form.

\paragraph{3. Argument normalization.}
The predicate \texttt{depth} has a three-argument signature
(object, reference-point, measure), but \texttt{mondial.kif}
contained hundreds of two-argument facts
\texttt{(depth Seas\_X 1234)}.
These were normalized to
\texttt{(depth Seas\_X SeaLevel (MeasureFn 1234 Meter))}.

\paragraph{4. Type confusion.}
In \texttt{emotion.kif}, \texttt{(contraryAttribute Pleasure Pain)}
treats \texttt{Pain} as an \texttt{Attribute}, but \SUMO{}'s own
taxonomy classifies \texttt{Pain} as a \texttt{PathologicalProcess}
(a subclass of \texttt{Process}, not \texttt{Attribute}).
This is a genuine ontological error: the axiom is well-formed in
\KIF{}'s free syntax but violates the intended sort discipline.

\paragraph{5. Semantic ambiguity.}
The predicate \texttt{monetaryValue} is declared with
\texttt{(domain monetaryValue 1 Physical)}, but the axioms defining
\texttt{WealthFn} use it for abstract financial quantities
(\texttt{PropertyFn}) that may not be \texttt{Physical} in the
intended sense.

\subsection{ATP Experience}

Pease and Sutcliffe~\cite{pease-sutcliffe-2007} reported sustained
efforts to make \SUMO{} amenable to \ATP{} systems at CASC.
The main challenges were:
\begin{itemize}[nosep]
  \item Row variable elimination (requires higher-order treatment or
    bounded unfolding).
  \item Self-referential axioms (\SUMO{} asserts things about its own
    predicates).
  \item Scale: first-order provers struggle with the full 20k-axiom
    set; modular decomposition is essential.
\end{itemize}

Benzm\"uller~\cite{benzmuller-faithful-embeddings} later showed how to
faithfully embed various logics (including \SUMO{}'s modal and
higher-order fragments) into classical \HOL{}, providing both ``deep''
(syntax-preserving) and ``shallow'' (semantics-preserving) embeddings.
This is closer to our approach but still treats axioms as opaque
assertions rather than extracting operational behavior.

%% ─────────────────────────────────────────────────────────────
\section{The Ontology Debugging Literature}
\label{sec:debugging}

\subsection{Glass-Box Debugging (Kalyanpur et al.~2005)}

Kalyanpur et al.~\cite{kalyanpur2005} introduced glass-box debugging
for \OWL{} ontologies, using the internals of description-logic tableaux
reasoners (specifically Pellet) to:
\begin{itemize}[nosep]
  \item Distinguish \emph{root} unsatisfiable classes from \emph{derived}
    ones (those that are unsatisfiable only because a root class is).
  \item Extract \emph{justifications}: minimal subsets of axioms
    responsible for an unsatisfiability.
  \item Present structured explanations to users in the Swoop
    ontology editor.
\end{itemize}

\textbf{Strength}: Pinpoints the \emph{root cause} of logical errors,
not just their symptoms.

\textbf{Limitation}: Restricted to description-logic (\OWL{}) ontologies.
\SUMO{}'s full \KIF{} (with row variables, self-application, and
unrestricted quantification) is far beyond \OWL{}'s expressiveness.
The tableaux-based approach does not extend to full \FOL{}.

\subsection{White-Box Testing (Keet \& \L{}awrynowicz 2017)}

Keet and \L{}awrynowicz~\cite{keet2017whitebox} proposed automatic test
generation from the syntactic structure of \FOL{} axioms:
\begin{itemize}[nosep]
  \item Each axiom pattern (universal, existential, disjunction, etc.)
    generates a family of positive and negative test cases.
  \item Tests are run against an \ATP{} to check whether the ontology
    entails (or correctly rejects) each test case.
  \item Redundant axioms are flagged when removing one does not change
    the test outcomes.
\end{itemize}

\textbf{Strength}: Fully automatic, works with \FOL{}.

\textbf{Limitation}: \FOL{} entailment is semi-decidable; the \ATP{}
may time out on hard test cases.
More fundamentally, the tests check \emph{what the ontology entails},
not \emph{whether those entailments match the intended domain
behavior}.

\subsection{Black-Box Testing (Keet \& \L{}awrynowicz 2017)}

The same authors proposed black-box testing using
WordNet~\cite{keet2017blackbox} as an external oracle:
\begin{itemize}[nosep]
  \item WordNet's synsets and hypernymy provide expected subsumption
    relationships.
  \item The ontology is queried: does it entail that $A$ subsumes $B$
    when WordNet says it should?
  \item Discrepancies identify potential errors or modeling choices.
\end{itemize}

\textbf{Strength}: External validation---not circular.

\textbf{Limitation}: Depends entirely on the quality of the
ontology-to-WordNet mapping.  WordNet's coverage is uneven across
domains, and mapping ambiguity introduces both false positives and
false negatives.

\subsection{OntoClean (Guarino \& Welty)}

OntoClean~\cite{guarino-welty-ontoclean} is a philosophical methodology
for evaluating taxonomic structures.
Each class is tagged with metaproperties:
\begin{itemize}[nosep]
  \item \textbf{Rigidity (+R/-R/$\sim$R)}: does the property hold
    necessarily for its instances?
  \item \textbf{Identity (+I/-I)}: does the class carry or inherit an
    identity criterion?
  \item \textbf{Unity (+U/-U/$\sim$U)}: do instances share a unifying
    whole-part relation?
  \item \textbf{Dependence (+D/-D)}: does the class's existence depend
    on another entity?
\end{itemize}

Constraint rules (e.g., ``an anti-rigid class must not subsume a rigid
one'') identify taxonomic errors.

\textbf{Strength}: Deep philosophical analysis of ``is this hierarchy
right?''---catching errors that no amount of logical consistency
checking can find.

\textbf{Limitation}: Entirely manual.  The metaproperties must be
assigned by a human expert, and there is no tool support for
propagation or conflict detection at scale.  We argue below
(\S\ref{sec:evidence-rigidity}) that the pipeline's evidence-valued
semantics can \emph{automate} rigidity assessment, replacing the
binary $\{0,1\}$ tags with graded evidence from corpus evaluation.

\subsection{Ontology Design Patterns (Gangemi \& Presutti)}

Gangemi and Presutti~\cite{gangemi-presutti-patterns} advocated
reusable ontology design patterns---modular templates for common
modeling situations (parthood, participation, classification, etc.).

\textbf{Strength}: Prevents errors by providing vetted templates.

\textbf{Limitation}: Post-hoc debugging is still needed for axioms
that do not fit a known pattern, and patterns do not address
behavioral correctness.

\subsection{Summary: A Layered but Disconnected Landscape}

No existing framework integrates all four correctness levels from
\S\ref{sec:correctness}.
Glass-box and white-box methods address levels 1--2 (syntax,
consistency).
OntoClean addresses level~3 (adequacy).
\emph{Level~4 (behavioral correctness) is unaddressed.}

Moreover, even where bugs are detected, the existing literature provides
limited guidance on \emph{repair}.  Axiom weakening~\cite{troquard2018weakening}
and optimal repair~\cite{baader2023pseudocontraction} offer principled repair
operators, but their theory is developed only for the lightweight description
logic $\mathcal{EL}$---far less expressive than \SUMO{}'s full \KIF{}.
Adimen-SUMO~\cite{alvez2012adimensumo} applied \ATP{}-based diagnosis to
re-engineer 88\% of \SUMO{}, but without behavioral or evidential signals
to disambiguate among repair candidates.

The pipeline we propose (\S\ref{sec:pipeline}) provides levels~1
and~4, and inherits level~2 from the underlying \ATP{} and type
system.  For level~3, the pipeline's evidence semantics can
\emph{partially automate} OntoClean's manual metaproperty
assignment (\S\ref{sec:evidence-rigidity}).
Crucially, it also provides a \emph{repair} methodology
(\S\ref{sec:repair}) that uses the pipeline's multi-level diagnostics
to guide principled corrections.

%% ─────────────────────────────────────────────────────────────
\section{The Pipeline: SUMO $\to$ GF $\to$ OSLF $\to$ WM}
\label{sec:pipeline}

\subsection{Overview}

\begin{center}
\begin{tabular}{@{}ccccccc@{}}
\SUMO{} & $\xrightarrow{\;\text{KIF parse}\;}$ &
\SUMO{}-\GF{} & $\xrightarrow{\;\text{OSLF synthesis}\;}$ &
\OSLF{} & $\xrightarrow{\;\text{evidence eval}\;}$ &
\WM{} \\[2pt]
\small axioms & &
\small GF abstract syntax & &
\small $\DiamondOp$, $\BoxOp$, native types & &
\small $\ev{n^+}{n^-}$ \\
\end{tabular}
\end{center}

Each stage catches a different class of errors:

\begin{center}
\small
\begin{tabular}{@{}lll@{}}
\toprule
\textbf{Stage} & \textbf{Bug class caught} & \textbf{Mechanism} \\
\midrule
1. \SUMO{} $\to$ \GF{} & Arity, sort errors & Grammar type discipline \\
2. \GF{} $\to$ \OSLF{} & Behavioral inconsistencies & $\DiamondOp$/$\BoxOp$ violation \\
3. \OSLF{} $\to$ \WM{} & Ungroundable assertions & Zero-evidence detection \\
& Scope ambiguities & Evidence ordering \\
\bottomrule
\end{tabular}
\end{center}

\subsection{Stage 1: SUMO $\to$ SUMO-GF}
\label{sec:stage1}

The Grammatical Framework (\GF{})~\cite{ranta2004gf} is a
grammar formalism with a type-theoretic abstract syntax and
language-specific concrete syntaxes (linearizations).
The key insight for ontology debugging is that \GF{}'s abstract
syntax is a \emph{typed} term language: every function signature
declares the sorts of its arguments and return value.

\paragraph{Mapping.}
\SUMO{} classes become \GF{} \emph{categories}, and \SUMO{} relations
become \GF{} \emph{function signatures}:

\begin{center}
\begin{tabular}{@{}ll@{}}
\toprule
\SUMO{} (\KIF{}) & \GF{} Abstract Syntax \\
\midrule
\texttt{(subclass Dog Animal)} & \texttt{cat Dog ; cat Animal} \\
\texttt{(domain bites 1 Animal)} & \texttt{fun bites : Animal -> Entity -> Prop ;} \\
\texttt{(domain bites 2 Entity)} & \\
\texttt{(instance Fido Dog)} & \texttt{fun Fido : Dog ;} \\
\bottomrule
\end{tabular}
\end{center}

\texttt{domain} and \texttt{range} declarations in \SUMO{} provide
the sort information; our existing \texttt{sumo-lean4} checker already
extracts these.
The translation is structurally identical to the HOL emitter in
\texttt{sumo-lean4/SumoKif/HolEmit.lean}, but targets \GF{} categories
instead of Lean types.

\paragraph{What this catches.}
An axiom like \texttt{(contraryAttribute Pleasure Pain)} requires
\texttt{Pain} to be of category \texttt{Attribute}.  \SUMO{}'s own
taxonomy declares \texttt{Pain} as a subclass of
\texttt{PathologicProcess}---a \texttt{Process}, not an
\texttt{Attribute}.  If the encoding faithfully follows \SUMO{}'s
declarations, the \GF{} type checker rejects the abstract
tree---\emph{before any reasoning occurs}.

Interestingly, the original \SUMO{}-\GF{} grammar~\cite{enache2012sumogf}
silently reclassifies \texttt{Pain} as \texttt{Ind~EmotionalState}
(an \texttt{Attribute}), making the axiom type-correct.  This is an
ontologically defensible choice---Pain is a sensation \emph{caused by}
a pathological process, not the process itself---but it represents an
undocumented repair.  Our pipeline
(Appendix~\ref{app:sumo-gf-experiment}) instead follows \SUMO{}'s
own declarations faithfully, detects the conflict, and logs the repair
decision with evidence and canary theorems
(\S\ref{sec:evidence-repair}).

\paragraph{Multilingual verbalization.}
As a bonus, \GF{}'s concrete syntax mechanism enables automatic
natural-language verbalization of \SUMO{} axioms, following Ranta
et al.~\cite{ranta-typeful-ontologies}.
An axiom that parses correctly in the grammar can be linearized to
English, Czech, or any other language with a concrete syntax---providing
a human-readable sanity check.

\subsection{Stage 2: GF $\to$ OSLF}
\label{sec:stage2}

Operational Semantics in Logical Form (\OSLF{}) takes a grammar
(presented as a \texttt{LanguageDef}: sorts, constructors, equations,
rewrite rules) and automatically synthesizes a behavioral type system.

\paragraph{LanguageDef.}
Our Mettapedia library represents a language as:
\begin{itemize}[nosep]
  \item \textbf{Sorts}: the categories (from \GF{}).
  \item \textbf{Terms} (\texttt{GrammarRule}): labeled constructors
    with typed parameters and syntax patterns.
  \item \textbf{Equations}: structural equality/normalization rules.
  \item \textbf{Rewrite rules}: one-step operational semantics.
  \item \textbf{Premises}: side conditions (freshness, congruence,
    relation queries).
\end{itemize}

The existing \GF{} $\to$ \OSLF{} bridge in Mettapedia
(\texttt{Languages/GF/OSLFBridge.lean}) translates 170 \GF{} RGL
function signatures into 985 grammar rules, and our pipeline
\texttt{langOSLF} synthesizes:
\begin{itemize}[nosep]
  \item A reduction relation $\Red$ from the rewrite rules.
  \item A possibility operator $\DiamondOp\,\phi$: ``a state satisfying
    $\phi$ is reachable.''
  \item A necessity operator $\BoxOp\,\phi$: ``all predecessors satisfy
    $\phi$.''
  \item A Galois connection $\DiamondOp \dashv \BoxOp$ (proven
    automatically from a generic theorem).
  \item Native types as (sort, predicate) pairs via the Grothendieck
    construction.
\end{itemize}

\paragraph{What this catches.}
Suppose \SUMO{} asserts ``every Human has a parent'' via:
\begin{verbatim}
  (=> (instance ?H Human)
      (exists (?P) (parent ?H ?P)))
\end{verbatim}
In the operational semantics derived from \SUMO{}'s axioms, if there
is no rewrite rule that produces a \texttt{parent} relation from a
\texttt{Human} state, then $s \not\models \DiamondOp(\texttt{parent}(h, p))$
for any state $s$ containing a Human~$h$.
This is a \emph{behavioral} bug: the ontology asserts an existential
that its own operational semantics cannot witness.

\subsection{Stage 3: OSLF $\to$ WM}
\label{sec:stage3}

The final stage interprets \OSLF{} behavioral types as evidence-valued
world-model queries.

\paragraph{Evidence.}
Following the PLN evidence quantale~\cite{goertzel2008pln},
evidence is a pair $\ev{n^+}{n^-} \in \mathbb{R}_{\geq 0}^{\infty}
\times \mathbb{R}_{\geq 0}^{\infty}$, where $n^+$ counts positive
evidence and $n^-$ counts negative evidence.
The strength (point estimate) is $s = n^+ / (n^+ + n^-)$.
Evidence forms a complete lattice and a commutative quantale under
coordinatewise multiplication.

\paragraph{World-model evaluation.}
Given a quantified formula $\phi$ (produced by the \GF{} $\to$ Store
$\to$ QFormula pipeline), the world-model evaluator $\mathrm{qsemE2}$
produces:
\[
  \mathrm{qsemE2}(\phi) \in \Ev
\]
where:
\begin{itemize}[nosep]
  \item Atomic formulas are evaluated by querying an evidence-atom
    semantics $I : \text{Pred} \times \text{Args} \times \text{Point}
    \to \Ev$.
  \item $\forall$ is interpreted as $\bigwedge$ (infimum over domain).
  \item $\exists$ is interpreted as $\bigvee$ (supremum over domain).
  \item $\DiamondOp$ and $\BoxOp$ are interpreted via reduction
    reachability.
\end{itemize}

\paragraph{What this catches.}

\begin{itemize}[nosep]
  \item \textbf{Ungroundable assertions}: if $\mathrm{qsemE2}(\phi) = \ev{0}{0}$
    (zero evidence), the ontology asserts something it cannot ground
    in any evidence source.
  \item \textbf{Contradictory evidence}: if $\mathrm{qsemE2}(\phi \wedge \psi)
    \ll \mathrm{qsemE2}(\phi) \otimes \mathrm{qsemE2}(\psi)$ (evidence
    for the conjunction is much less than the product of evidence for
    the conjuncts), the axioms are in tension.
  \item \textbf{Scope ambiguities}: for quantified \SUMO{} axioms, the
    inverse-scope reading may produce different evidence than the
    surface reading.  Our \texttt{scope\_ordering\_NT} theorem in
    Mettapedia proves that inverse $\leq$ surface in the evidence
    lattice, so a significant gap flags a scope-sensitive assertion.
\end{itemize}

\subsection{Compositionality}

A key advantage of the pipeline is \emph{compositionality}.
\GF{}'s abstract syntax is Frege-compositional by construction:
the meaning of a complex expression is determined by the meanings
of its parts and their mode of combination.
\OSLF{}'s type synthesis preserves this: the behavioral type of a
compound term is determined by the behavioral types of its subterms.
The world-model evaluator is a compositional denotational semantics.

This means that if two axioms are individually well-typed and
behaviorally correct, their combination is type-correct under the
compositional typing theorem used in this pipeline slice.
\emph{Compositional bugs}---where two individually valid axioms
interact to produce an ill-typed behavior---are caught at the \OSLF{}
level by the native type discipline, specifically by the Grothendieck
construction over the subobject fibration
(\texttt{CategoryTheory/NativeTypeTheory.lean}).

\subsection{Automating OntoClean via Evidence-Valued Rigidity}
\label{sec:evidence-rigidity}

OntoClean's central metaproperty, \emph{rigidity}, asks: ``if $x$ is
a $C$, must $x$ \emph{necessarily} be a $C$?''  Traditional OntoClean
assigns a binary tag: $+\mathrm{R}$ (rigid), $-\mathrm{R}$
(non-rigid), or ${\sim}\mathrm{R}$ (anti-rigid).  The assignment is
manual and requires philosophical expertise.

We observe that the pipeline's evidence-valued semantics provides a
natural \emph{automation} of rigidity assessment, replacing the
binary $\{0,1\}$ tag with a graded measure in $[0,1]$.

\paragraph{Formalization.}
For each ontology class $C$, define the \emph{rigidity evidence} as
the world-model evaluation of the claim that some instance of $C$
could cease to be a $C$:
\[
  \mathrm{rig\text{-}ev}(C) \;:=\;
    \mathrm{qsemE2}\bigl(\exists x.\;
      \mathrm{instanceOf}(x, C) \;\wedge\;
      \DiamondOp(\lnot\,\mathrm{instanceOf}(x, C))\bigr)
\]
Here $\DiamondOp$ is the pipeline's behavioral possibility operator,
which checks whether $\lnot\,\mathrm{instanceOf}(x, C)$ is
reachable under the ontology's rewrite dynamics.  The result is an
evidence pair $\ev{n^+}{n^-}$.

The \emph{graded rigidity} of $C$ is then:
\[
  \rho(C) \;:=\; 1 - s(\mathrm{rig\text{-}ev}(C))
  \;=\; 1 - \frac{n^+}{n^+ + n^-}
\]
where $s$ is the PLN strength.  When no evidence is found for
instances losing class membership ($n^+ = 0$), $\rho(C) = 1$
(fully rigid).  When abundant evidence exists ($n^+ \gg n^-$),
$\rho(C) \approx 0$ (anti-rigid).

\paragraph{Examples.}
Consider three \SUMO{} classes evaluated against a commonsense
corpus:

\begin{center}
\small
\begin{tabular}{@{}llcl@{}}
\toprule
\textbf{Class} & \textbf{Evidence for change} &
  $\boldsymbol{\rho(C)}$ & \textbf{OntoClean tag} \\
\midrule
\texttt{Human} & No factual examples of becoming non-Human
  & $\approx 1.0$ & $+$R \\
\texttt{Student} & Graduation, dropout ($n^+ \gg 0$)
  & $\approx 0.1$ & ${\sim}$R \\
\texttt{Physical} & No evidence of physical$\to$abstract
  & $= 1.0$ & $+$R \\
\texttt{President} & Term endings, impeachment ($n^+ > 0$)
  & $\approx 0.2$ & ${\sim}$R \\
\bottomrule
\end{tabular}
\end{center}

\paragraph{Domain-relative rigidity.}
A key advantage of the evidence-based approach is that rigidity
becomes \emph{domain-relative}: the same class can have different
rigidity values depending on the world model's corpus.  A medical
ontology evaluated against clinical records will assign $\rho(\texttt{Human})
\approx 1.0$, while a science-fiction ontology might assign a lower
value.  This is not a defect but a feature: rigidity genuinely
\emph{should} depend on the domain of discourse.  Traditional
OntoClean obscures this by forcing a single universal tag.

\paragraph{Constraint propagation.}
OntoClean's constraint rules translate directly to evidence
inequalities.  The rule ``an anti-rigid class must not subsume a
rigid one'' becomes: if $C_1 \sqsubseteq C_2$ and
$\rho(C_2) > \rho(C_1)$, then the subsumption is suspect---a child
class is more changeable than its parent, violating monotonicity of
rigidity along the taxonomy.  These inequalities can be checked
\emph{automatically} across the entire ontology once graded rigidity
values are computed.

\paragraph{Relation to other metaproperties.}
The same evidence-valued approach extends to OntoClean's other
metaproperties.  \emph{Identity} can be assessed by evaluating
evidence for identity criteria: $\mathrm{qsemE2}(\exists f.\;
\forall x,y : C.\; f(x) = f(y) \to x = y)$.  \emph{Unity} can be
assessed via evidence for part-whole coherence.  We leave the formal
development of these extensions to future work, noting that the
pipeline's evidence quantale provides the necessary algebraic
structure in each case.

%% ─────────────────────────────────────────────────────────────
\section{From Detection to Repair: The Formalization-Driven Loop}
\label{sec:repair}

The pipeline of \S\ref{sec:pipeline} \emph{detects} bugs across four
correctness levels.  But detection is only half the story: a principled
\emph{repair} methodology is needed to close the loop.  This section
surveys existing ontology repair approaches, identifies their limitations
for full-\FOL{} ontologies like \SUMO{}, and proposes a convergent
formalization-driven repair loop that exploits the pipeline's behavioral
and evidential signals.

\subsection{The Ontology Repair Literature}

\paragraph{Justification-based repair.}
Kalyanpur et al.~\cite{kalyanpur2005} introduced \emph{Minimal
Unsatisfiability-Preserving Sub-TBoxes} (MUPS): the smallest subset of
axioms sufficient for a concept to be unsatisfiable.  Horridge et
al.~\cite{horridge2008laconic} refined this to \emph{laconic
justifications} that eliminate redundancy within justifications.
Repair proceeds by removing or modifying at least one axiom from each
justification.

\textbf{Limitation}: These methods require a complete reasoner (typically
a DL tableaux); they do not extend to \SUMO{}'s full \KIF{} with row
variables and self-application.  Moreover, deletion-based repair loses
information.

\paragraph{Axiom weakening.}
Troquard et al.~\cite{troquard2018weakening} proposed \emph{weakening}
axioms rather than deleting them, using refinement operators from
inductive logic programming.  For example,
$A \sqsubseteq B \sqcap C$ may be weakened to $A \sqsubseteq B$
(dropping a conjunct) rather than deleted entirely.  This preserves
significantly more of the original knowledge.

Baader and Kriegel~\cite{baader2023pseudocontraction,baader2022fixedpremise}
showed that optimal repair in $\mathcal{EL}$ TBoxes corresponds
precisely to \emph{pseudo-contraction} in AGM belief change theory:
the repair that removes the fewest entailments while restoring
consistency.  They also introduced \emph{fixed-premise repairs}, where
some axioms are trusted (``fixed'') and only the remainder is modified.

\textbf{Limitation}: The theory is developed only for $\mathcal{EL}$,
with recent extensions to Horn-DLs~\cite{baader2022fixedpremise}.
\SUMO{}'s full first-order expressiveness (row variables, unrestricted
quantification, self-reference) is far beyond these fragments.

\paragraph{Adimen-SUMO.}
Alvez et al.~\cite{alvez2012adimensumo} applied \ATP{}-based diagnosis
directly to \SUMO{}: for each axiom~$\alpha$, attempt to prove
$\neg\alpha$ from the remaining axioms.  Success indicates that $\alpha$
contradicts the rest of the knowledge base.  This methodology
re-engineered approximately 88\% of \SUMO{}, finding incorrectly
defined axioms, redundancies, and non-desirable properties.

\textbf{Limitation}: Catches logical inconsistency but not type errors,
behavioral bugs, or evidential gaps.  The methodology is purely
deductive and provides no guidance on \emph{which} repair to choose
when multiple corrections are possible.

\paragraph{SUMO-GF type checking.}
Enache~\cite{enache2012sumogf} translated 17 \SUMO{} modules
to \GF{} abstract syntax, using Martin-L\"of dependent types to encode
\SUMO{}'s sort discipline.  Of the quantified axioms, \textbf{64\%
passed} type-checking; the remaining 36\% broke down as: 12\%
rejected by the type checker (genuine sort/arity errors), 12\%
inexpressible because variable types are inferred from usage context,
11\% using predicates not translated to \GF{} (\texttt{range},
\texttt{domain}, \texttt{subclass}), and 1\% other.  This is the
direct predecessor of our Stage~1, and the most relevant prior art.

Crucially, examination of the original \SUMO{}-\GF{} source
(\texttt{Merge.gf}, \texttt{MidLevelOntology.gf}, distributed with
GF~3.3) reveals that the encoding contains \emph{silent repairs}:
classification choices that differ from \SUMO{}'s own taxonomy but
are not documented as corrections.  For example, \SUMO{} declares
\texttt{Pain} as a subclass of \texttt{PathologicProcess} (a
\texttt{Process}), but the \GF{} grammar declares
\texttt{fun Pain : Ind EmotionalState} (an \texttt{Attribute}).
This reclassification makes \texttt{contraryAttribute(Pleasure, Pain)}
type-correct---a defensible ontological judgment (Pain is a sensation
caused by a pathological process, not the process itself)---but it
was made implicitly, with no audit trail or canary theorem to verify
correctness.

\textbf{Limitation}: Type errors were detected but not automatically
repaired; where repairs were made, they were undocumented.  The project
provides no systematic methodology for distinguishing genuine \SUMO{}
bugs from encoding artifacts, nor for verifying that silent repairs
preserve the intended semantics.

\paragraph{The gap.}
No existing system provides a \emph{closed-loop} repair methodology
for full-\FOL{} ontologies that (1)~detects errors across multiple
levels, (2)~uses structured signals to \emph{disambiguate} among repair
candidates, (3)~verifies repairs via proven invariants, and
(4)~converges.  Our proposal addresses this gap.

\subsection{The Convergent Repair Loop}
\label{sec:repair-loop}

We propose a six-phase loop that proceeds top-down through \SUMO{}'s
class hierarchy, starting from \texttt{Entity}:

\begin{enumerate}
  \item \textbf{Formalize.}  Translate the next layer of \SUMO{} axioms
    to \GF{} abstract syntax (Stage~1 of the pipeline).  Proceed
    top-down: first \texttt{Entity} and its immediate subclasses, then
    \texttt{Physical}/\texttt{Abstract}, and so on.

  \item \textbf{Detect.}  Run the pipeline diagnostics:
    \begin{itemize}[nosep]
      \item \GF{} type checker: arity and sort errors.
      \item \OSLF{} synthesis: behavioral type violations
        ($\DiamondOp/\BoxOp$ ill-formedness).
      \item \WM{} evaluator: zero-evidence assertions, scope
        ambiguities.
    \end{itemize}

  \item \textbf{Disambiguate.}  For each detected error, the pipeline
    provides structured context that narrows the space of candidate
    repairs:
    \begin{itemize}[nosep]
      \item \emph{GF type context}: bidirectional type inference gives
        the expected sort at the error site (e.g., ``expected
        \texttt{Attribute}, got \texttt{Process}''), directly suggesting
        which argument or declaration is wrong.
      \item \emph{OSLF behavioral types}: different repairs produce
        different operational semantics; the one preserving intended
        reduction behavior is preferred.
      \item \emph{WM evidence scores}: repairs can be ranked by
        evidential groundability---the fix best supported by evidence
        is preferred.
    \end{itemize}

  \item \textbf{Repair.}  Propose a minimal \KIF{} patch.
    Following the axiom-weakening
    principle~\cite{troquard2018weakening}, weakening is preferred over
    deletion: strengthen a sort constraint, split a concept, or add a
    missing argument---rather than removing an axiom entirely.

  \item \textbf{Verify.}  Prove a \emph{canary theorem}
    (Definition~\ref{def:canary}) that encodes the intended consequence
    of the repaired axiom.  If the canary fails, the repair is
    incorrect and the loop returns to step~3.

  \item \textbf{Continue.}  Add the canary to the growing regression
    suite and proceed to the next layer.
\end{enumerate}

\begin{principle}[Convergence]
The loop converges because: (a)~each \SUMO{} layer has finitely many
axioms; (b)~each successful repair strictly reduces the error count at
that layer; (c)~canary theorems prevent regression (a repair at layer
$k$ that breaks a canary at layer $j < k$ is immediately caught).
The canary suite grows monotonically, providing an ever-stronger
invariant.
\end{principle}

\begin{remark}[Connection to autoformalization]
The repair loop has the same structure as the LLM autoformalization
loop~\cite{autoformalization-survey}: formalize $\to$ check $\to$
error feedback $\to$ retry.  The key difference is the \emph{richness
of the error signal}: where autoformalization typically receives only
a type error message, our pipeline provides type context, behavioral
diagnostics, and evidence scores---three independent disambiguation
channels.
\end{remark}

\subsection{Canary Theorems}
\label{sec:canary}

\begin{definition}[Canary theorem]
\label{def:canary}
A \emph{canary theorem} for a repair $r$ applied to axiom $\alpha$ is
a small, proven proposition $\gamma_r$ such that:
\begin{enumerate}[nosep]
  \item $\gamma_r$ is a direct intended consequence of the repaired
    axiom $\alpha[r]$.
  \item $\gamma_r$ is \emph{falsified} by any incorrect repair $r'$
    (i.e., $\alpha[r'] \not\models \gamma_r$).
  \item $\gamma_r$ is proven in the proof assistant (Lean), not merely
    tested via \ATP{} query.
\end{enumerate}
\end{definition}

\begin{example}[Classmate arity canary]
After repairing \texttt{classmate} from
\texttt{IrreflexiveRelation} membership to an explicit ternary
irreflexivity axiom (\S\ref{sec:example}), the canary is:
\[
  \gamma_{\texttt{classmate}} :\;
  \forall\, a\, c.\; \neg\,\texttt{classmate}(a, a, c)
\]
This is provable from the repaired axiom.  If a future edit
accidentally changes the arity (e.g., to quaternary), the canary
immediately fails.
\end{example}

\begin{example}[Pain/Attribute canary]
\label{ex:pain-canary}
After splitting \texttt{Pain} into \texttt{PainProcess} (a subclass
of \texttt{Process}) and \texttt{PainSensation} (an \texttt{Attribute}),
the canary suite includes:
\begin{align*}
  \gamma_1 &:\; \texttt{instance}(\texttt{PainProcess}, \texttt{Process}) \\
  \gamma_2 &:\; \texttt{instance}(\texttt{PainSensation}, \texttt{Attribute}) \\
  \gamma_3 &:\; \texttt{contraryAttribute}(\texttt{Pleasure},
    \texttt{PainSensation}) \text{ is well-typed} \\
  \gamma_4 &:\; \texttt{PainProcess} \neq \texttt{PainSensation}
\end{align*}
Any repair that merges them back or assigns the wrong metaclass breaks
at least one canary.
\end{example}

Canary theorems relate to Keet and \L{}awrynowicz's white-box
testing~\cite{keet2017whitebox}: both generate test obligations from
axiom structure.  The key difference is that canaries are \emph{proven}
in a proof assistant, not queried via a semi-decidable \ATP{} call.
A canary either holds or it does not---there is no timeout ambiguity.

\subsection{Disambiguation via Behavioral Evidence}
\label{sec:disambiguation}

The central advantage of the pipeline for repair---as opposed to purely
type-directed or purely deductive approaches---is that it provides
\emph{multiple independent signals} for choosing among repair
candidates.  We illustrate with the Pleasure/Pain bug
(\S\ref{sec:sumo}).

\paragraph{The ambiguity.}
The type error is clear: \texttt{contraryAttribute(Pleasure, Pain)}
requires \texttt{Pain} to be an \texttt{Attribute}, but \SUMO{}'s
taxonomy says \texttt{Pain} is a \texttt{Process}.
Three repairs are possible:
\begin{enumerate}[label=(\alph*),nosep]
  \item Reclassify \texttt{Pain} as \texttt{Attribute} (breaking its
    process axioms).
  \item Generalize \texttt{contraryAttribute} to accept
    \texttt{Process} (weakening the sort constraint).
  \item Split: create \texttt{PainSensation} $:$ \texttt{Attribute}
    alongside the existing \texttt{Pain} $:$ \texttt{Process}
    (concept refinement).
\end{enumerate}

\paragraph{Type signal alone.}
A pure type checker (like \SUMO{}-\GF{}~\cite{enache2012sumogf}) reports
the mismatch but cannot distinguish (a), (b), and (c).  All three
restore type-correctness.

\paragraph{Behavioral signal.}
\OSLF{} synthesis reveals that \texttt{Pain} participates in reduction
rules (it has an operational semantics as a process: onset, duration,
intensification, subsidence).  Repair~(a) would orphan these rules,
producing $\BoxOp$-violations for any scenario involving pain dynamics.
Repair~(b) would allow \texttt{contraryAttribute} to compare processes
and attributes, producing ill-typed behavioral consequences.
Repair~(c) preserves all existing reduction rules for
\texttt{Pain}/\texttt{PainProcess} while creating a new, behaviorally
inert \texttt{PainSensation} attribute.
\emph{The behavioral signal selects~(c).}

\paragraph{Evidence signal.}
The \WM{} evaluator can further rank candidates: medical and
neuroscience corpora provide strong evidence for Pain-as-Process
($\ev{n^+}{n^-}$ with high $n^+$), while folk-psychological usage
supports Pain-as-Attribute-like quality.  The split repair~(c)
maximizes total evidence across both readings.

\begin{table}[ht]
\centering
\caption{Repair disambiguation for the Pain/Attribute bug.}
\label{tab:disambiguation}
\small
\begin{tabular}{@{}lccc@{}}
\toprule
\textbf{Signal} & \textbf{(a) Reclassify} & \textbf{(b) Generalize} & \textbf{(c) Split} \\
\midrule
Type-correct? & \checkmark & \checkmark & \checkmark \\
Behavioral coherence & \texttimes & \texttimes & \checkmark \\
Evidence grounding & partial & partial & maximal \\
Canary suite passes & $\gamma_{1,3}$ fail & $\gamma_3$ fails & all pass \\
\bottomrule
\end{tabular}
\end{table}

This multi-signal disambiguation is what distinguishes the pipeline
approach from all prior ontology repair methods.  Type-directed repair
(SUMO-GF, LeanSUMO) narrows the space; behavioral types eliminate
operationally incoherent candidates; evidence scores rank the remainder.

%% ─────────────────────────────────────────────────────────────
\section{The Evidence Model: Structured Assurance for Repair Decisions}
\label{sec:evidence-model}

Detection alone does not justify a repair.  The gap between
``\texttt{.unsat} detected'' and ``repair safely applicable'' requires
structured evidence, multi-agent review, and quantitative assurance
scoring.  This section describes the \emph{Evidence Model}---a formal
framework implemented in Lean~4 (\texttt{EvidenceModel.lean},
282~lines)---that governs which repairs may be applied and when.

\subsection{Repair Archetypes and Beta Priors}

Not all repairs carry the same epistemic risk.  A syntax fix
(\texttt{exist}$\to$\texttt{exists}) is nearly certain; reclassifying
a concept across ontological branches requires deep judgment.
We define six \emph{repair archetypes}, each with a Beta-distribution
prior $(\alpha_0, \beta_0)$ encoding how much evidence is needed
before acting:

\begin{center}\small
\begin{tabular}{@{}lccl@{}}
\toprule
\textbf{Archetype} & $\alpha_0$ & $\beta_0$ & \textbf{Example} \\
\midrule
\texttt{syntaxArity} & 12.0 & 1.0 & \texttt{exist}$\to$\texttt{exists}, missing args \\
\texttt{argSwapTyped} & 10.0 & 1.5 & D24: agent arguments swapped \\
\texttt{domainNarrowing} & 7.0 & 3.0 & D26: Group$\to$AutonomousAgent \\
\texttt{hierarchyGap} & 6.0 & 3.5 & D23: Artifact bypasses SelfConnectedObject \\
\texttt{ontologyReclassification} & 4.0 & 4.0 & D1: Pain Process$\to$Attribute \\
\texttt{worldSemanticClaim} & 3.0 & 5.0 & Deep semantic/world-model judgment \\
\bottomrule
\end{tabular}
\end{center}

Higher $\alpha_0$ relative to $\beta_0$ means the archetype is
easier to trust from less evidence.  Syntax fixes have prior mean
$12/13 \approx 0.92$; world-semantic claims have prior mean
$3/8 = 0.375$---starting from skepticism.

\subsection{Evidence Atoms and Source Weighting}

Every piece of evidence---from a proven-sound \texttt{checkLang}
result to a literature citation to an LLM review vote---is
normalized into a common \emph{evidence atom} schema carrying:
claim ID, decision ID, archetype, direction (supports/contradicts),
source type, formal expression, strength, confidence, independence
group (to avoid double-counting correlated evidence), and a
reproducibility hash.

Source types are weighted by reliability:

\begin{center}\small
\begin{tabular}{@{}lrl@{}}
\toprule
\textbf{Source type} & \textbf{Weight} & \textbf{Rationale} \\
\midrule
\texttt{checkLang} (proven sound) & 1.00 & Verified by \texttt{checkLangUsing\_sat\_sound} \\
\texttt{stratumScan} & 0.75 & \texttt{constraint\_check.py} output \\
\texttt{kifCitation} & 0.65 & Direct KIF file:line reference \\
\texttt{engGFParsed} & 0.60 & GF runtime: unique parse tree + roundtrip verified \\
\texttt{humanReview} & 0.60 & Human expert vote \\
\texttt{kifRendered} & 0.50 & Deterministic KIF $\to$ controlled English template \\
\texttt{docGFParsed} & 0.50 & Doc definition: GF-parsed + KIF cross-checked \\
\texttt{literature} & 0.45 & Encyclopedia, medical standard, etc. \\
\texttt{llmReview} & 0.35 & Claude Code / Codex / GPT-5.2 Pro vote \\
\texttt{englishDoc} & 0.30 & Corroborative hint from \texttt{(documentation ...)} \\
\bottomrule
\end{tabular}
\end{center}

The \texttt{checkLang} source type has maximal weight because its
results are \emph{proven sound}: Theorem
\texttt{checkLangUsing\_sat\_sound} guarantees that a \texttt{.sat}
result implies the semantic property actually holds.
The three GF-derived sources (\texttt{kifRendered}, \texttt{engGFParsed},
\texttt{docGFParsed}) form an ascending trust hierarchy:
template rendering (deterministic but pattern-limited),
GF runtime parsing (typed grammar with AST equivalence verification),
and documentation cross-checking (GF-parsed definitions verified
against KIF facts).
LLM reviews have low weight due to correlation (multiple LLMs trained
on overlapping data).  \texttt{englishDoc} is the weakest source:
prose-level hints from SUMO documentation strings, used only as
corroborative evidence and never as sole authority.

\subsection{Assurance Scoring}

Evidence atoms contribute weighted pseudo-counts to the archetype's
Beta prior.  A supporting atom from source~$s$ with strength~$t$
and confidence~$c$ adds $w_s \cdot t \cdot c$ to $\alpha$; a
contradicting atom adds the same to $\beta$.  The
\emph{assurance number} (``ass-number'') is the lower bound of
the 95\% credible interval of $\text{Beta}(\alpha, \beta)$:
\[
  \text{ass} = F^{-1}_{\text{Beta}(\alpha,\beta)}(0.025)
\]
where $F^{-1}$ is the inverse CDF\@.  This is a conservative
estimate: we are 97.5\% confident the true ``goodness'' of the
repair exceeds $\text{ass}$.

\paragraph{Thresholds.}
\begin{itemize}[nosep]
  \item $\text{ass} \geq 0.98$: eligible for auto-application
    (if review quorum met)
  \item $0.93 \leq \text{ass} < 0.98$: hold for additional reviewer
  \item $\text{ass} < 0.93$: remain logged/deferred
\end{itemize}

\subsection{Multi-Agent Review Protocol}

Classification decisions for full-\FOL{} ontologies benefit from
diverse perspectives.  The Evidence Model supports a multi-agent
review protocol where each reviewer (human or AI) submits a
\emph{review vote}: decision ID, reviewer identity, verdict
(approve/reject/abstain), confidence, and hashes of the rationale
and evidence bundle reviewed.

Planned reviewers for the pipeline (review votes not yet
populated; all 26 decisions are at \texttt{pendingReview}):
\begin{itemize}[nosep]
  \item \textbf{Claude Code}: primary development agent, produces
    initial evidence and repair proposals
  \item \textbf{Codex}: independent review with deep context;
    identified D20/D21 status corrections and evidence policy gaps
  \item \textbf{GPT-5.2 Pro}: cross-validation with different
    training distribution
  \item \textbf{Human}: final authority, especially for
    world-semantic judgments
\end{itemize}

LLM reviewers in the same independence group (e.g., all GPT-family
models) have their evidence down-weighted to avoid inflating
confidence from correlated sources.

\paragraph{Review lifecycle.}
Each repair decision transitions through states:
$\texttt{pendingReview} \to \texttt{reviewApproved} \to
\texttt{applied} \to \texttt{verified}$ (or $\to \texttt{rejected}$
if strong contradictory evidence emerges).  No KIF or SUMO-GF
repairs are applied until the decision reaches
\texttt{reviewApproved}.

\subsection{WM Witnesses: Persisting checkLang Results}

The comprehensive audit (\S\ref{sec:ntt-extraction}) produces
\texttt{checkLang} results at build time.  The Evidence Model
persists these as structured \emph{WM witnesses} in the repair
log, each recording: witness ID, query, expected result,
observed result, pass/fail, and source reference.

For decisions D22--D26, the current witness inventory is:

\begin{center}\small
\begin{tabular}{@{}llllll@{}}
\toprule
\textbf{ID} & \textbf{Decision} & \textbf{Query (reachabilityTest)} & \textbf{Exp.} & \textbf{Obs.} & \textbf{Pass} \\
\midrule
D22 & LinguisticExpression$\not\to$Object & LE $\to$ object & unsat & unsat & \checkmark \\
D23 & Artifact$\not\to$SelfConnObj & Artifact $\to$ selfconnobj & unsat & unsat & \checkmark \\
D24a & CogAgent$\not\to$Process (pre) & CA $\to$ process & unsat & unsat & \checkmark \\
D24b & CogAgent$\to$Object (post) & CA $\to$ object & sat & sat & \checkmark \\
D26a & CogAgent$\not\to$Group (pre) & CA $\to$ group & unsat & unsat & \checkmark \\
D26b & Group$\to$Object (positive) & Group $\to$ object & sat & sat & \checkmark \\
\bottomrule
\end{tabular}
\end{center}

All 6 witnesses pass.  Combined with the 14 existing testable
decisions from D1--D20, the comprehensive audit now covers
\textbf{20 passed / 0 failed / 8 N/A} across 26 decisions.
Artifact pointers (snapshot): \texttt{Mettapedia/Languages/GF/SUMO/data/wm\_witnesses.jsonl},
\texttt{Mettapedia/Languages/GF/SUMO/data/decision\_coverage.json},
\texttt{Mettapedia/Languages/GF/SUMO/export\_wm\_witnesses.py}.

\subsection{Evidence-Quantale Compatibility}
\label{sec:ev-quantale}

The Evidence Model's $(n^+, n^-)$ pseudo-counts align directly with
\texttt{Mettapedia.Logic.EvidenceQuantale}, which models evidence as
pairs in $\mathbb{R}_{\geq 0}^{\infty} \times
\mathbb{R}_{\geq 0}^{\infty}$ forming a quantale under
component-wise addition.  The conversions \texttt{toSTV}/\texttt{ofSTV}
bridge to PLN's \texttt{SimpleTruthValue} (strength, confidence).
This means the repair pipeline's evidence atoms can be consumed by
PLN inference---e.g., to propagate evidence about a concept's
classification to axioms that use it.

\paragraph{Policy: no repairs without evidence.}
The critical policy decision is: \texttt{autoFixable} marks a
candidate \emph{eligible} for auto-application, not ``apply now.''
Application requires: (1)~sufficient evidence atoms across
required witness types (pre-fix, post-fix, regression),
(2)~assurance number $\geq$ threshold, and (3)~review quorum met.
This prevents premature repairs from propagating errors deeper into
the ontology.

\subsection{GF Runtime Evidence Pipeline}
\label{sec:gf-evidence}

Beyond proven-sound \texttt{checkLang} witnesses and manual KIF
citations, we construct evidence atoms \emph{automatically} from
SUMO's own axioms and documentation using the Grammatical
Framework~(GF) as a typed grammar runtime.  This yields three
additional evidence lanes, each with distinct trust characteristics.

\paragraph{Lane~A: Template-rendered controlled English (\texttt{kifRendered}).}
A deterministic template renderer converts each KIF axiom into
controlled English matching GF linearization patterns:
\texttt{(subclass~X~Y)} $\mapsto$ ``\texttt{X~is~a~subclass~of~Y}'',
\texttt{(instance~X~Y)} $\mapsto$ ``\texttt{X~is~an~instance~of~Y}'',
and so on through 22~relation patterns covering quantifiers
(\texttt{forall}, \texttt{exists}), connectives
(\texttt{if\ldots then}, \texttt{iff}, \texttt{and}, \texttt{or},
\texttt{not}), and thematic roles (\texttt{agent}, \texttt{patient},
\texttt{instrument}).  Running this over the 937 repair-relevant axioms
produces 10{,}584~template renderings, all with 100\% AST roundtrip
fidelity (the rendered English, when parsed back through the template
grammar, reconstructs the original KIF AST).  This lane contributes
805~evidence atoms at weight~0.50.

\paragraph{Lane~B: GF runtime parsing (\texttt{engGFParsed}).}
A minimal RGL-free GF grammar (\texttt{AxiomCE.gf}) encodes
the controlled English patterns as typed abstract syntax:
\begin{verbatim}
  cat Stmt ; Formula ; Term ; ArgN ;
  fun SubclassStmt  : Term -> Term -> Stmt ;
      InstanceStmt  : Term -> Term -> Stmt ;
      ExistsFml     : Term -> Formula -> Formula ;
      ImplFml       : Formula -> Formula -> Formula ;
      ...
\end{verbatim}
with an auto-generated closed lexicon of 17{,}718~terms
(all identifiers appearing in SUMO axioms and documentation
cross-references).  The grammar compiles to a PGF
(Portable Grammar Format) file and runs as a GF subprocess.

For each axiom, the pipeline: (1)~normalizes the template-rendered
English (CamelCase $\to$ underscore, comma spacing);
(2)~parses through GF (\texttt{p~-cat=Stmt});
(3)~linearizes the parse tree back (\texttt{l}); and
(4)~verifies \emph{AST equivalence}: the GF tree's top-level
constructor must match the KIF head relation
(e.g., \texttt{SubclassStmt} for \texttt{(subclass~\ldots)}).
Only axioms passing all three gates---unique parse, string
roundtrip, and AST equivalence---produce \texttt{engGFParsed}
atoms.

\medskip\noindent\textbf{Disambiguation.}
A key engineering decision is the separation of
\emph{canonical} (machine) English from \emph{human-friendly}
English.  The canonical form parenthesizes every compound
sub-formula: \texttt{if~(~A~and~(~B~and~C~)~)~,~then~(~D~)}
--- with a grammar constructor \texttt{ParenFml} producing
\texttt{``(''~++~f.s~++~``)''}.  $N$-ary conjunction
right-associates: \texttt{A~and~(~B~and~C~)} rather than
the ambiguous \texttt{A~and~B~and~C}.  This eliminates
all parse ambiguity: the GF chart parser returns exactly
one tree for every parenthesized axiom.

\medskip\noindent\textbf{Results.}
Of 937~repair-relevant axioms: 802~(85.6\%) parse uniquely with
roundtrip and AST match; 0~parse ambiguously; 135~fail to parse.
The 135~failures reflect missing grammar coverage for
SUMO-specific constructs (checkpoint 2026-02-23):

\begin{center}\small
\begin{tabular}{@{}lr@{}}
\toprule
\textbf{Gap family} & \textbf{Count} \\
\midrule
Complex connectives (nested if/for-all) & 47 \\
Multi-variable existentials ($\exists$2) & 20 \\
Multi-variable existentials ($\exists$3) & 16 \\
Statement reification (\texttt{KappaFn}) & 16 \\
Multi-variable existentials ($\exists$4+) & 15 \\
Spatial verbs (\texttt{ends\_up\_at}, etc.) & 6 \\
Row-variable predicates (\texttt{@ROW}) & 5 \\
IFF at top level (\texttt{iff\_toplevel}) & 5 \\
Obliged-perform constructs & 3 \\
Other & 2 \\
\bottomrule
\end{tabular}
\end{center}

These are grammar coverage gaps, not GF expressiveness
limits.  Grammar improvements between v1 and this checkpoint
resolved: functional term equality (\texttt{equal}→``is equal
to''), 16 new verb/relation patterns, and formula-argument
parenthesization for higher-order relations (\texttt{holdsDuring},
\texttt{believes}, \texttt{desires}).  The aggregate count
reflects SUMO's large vocabulary of specialized relations.
This lane contributes 514~evidence atoms at weight~0.60.
Artifact pointers (snapshot): \texttt{Mettapedia/Languages/GF/SUMO/data/axioms\_gf\_parse.jsonl},
\texttt{Mettapedia/Languages/GF/SUMO/data/gf\_parse\_provenance.json},
\texttt{Mettapedia/Languages/GF/SUMO/data/gf\_parsed\_evidence\_atoms.jsonl},
\texttt{Mettapedia/Languages/GF/SUMO/sumo\_axiom\_gf\_parse.py}.

\paragraph{Lane~C: Documentation parsing (\texttt{docGFParsed}).}
SUMO's \texttt{(documentation~Term~EnglishLanguage~``\ldots'')}
entries contain encyclopedia-style definitions.  We extract
definitional head claims using four regex patterns:
\begin{enumerate}[nosep]
  \item ``A~\$X is a~\$Y'' or ``\$X is a~\$Y'' $\to$ subclass claim
  \item ``\$X is an instance of~\$Y'' $\to$ instance claim
  \item ``Any~\$X that\ldots'' $\to$ implied subclass of~X's parent
  \item ``A~\$X that\ldots'' $\to$ typed particular (implied subclass)
\end{enumerate}

Extracted heads are normalized and parsed through the same
\texttt{AxiomCE} grammar.  A KIF cross-check classifies each
parsed claim as \emph{consistent} (claim matches a KIF
\texttt{subclass}/\texttt{instance} fact), \emph{contradicts}
(KIF has a different superclass for the same subject), or
\emph{no\_kif\_data} (no matching KIF fact found).

\medskip\noindent\textbf{Results.}
Of 13{,}832~documentation entries: 2{,}130~yield extractable
head claims; 1{,}555~of these parse uniquely through GF (73\%).
KIF cross-checking: 425~consistent, 225~contradictory,
905~no KIF data.  This lane contributes 60~evidence atoms at
weight~0.50, covering 8~repair decisions.  Contradictory atoms
(where documentation disagrees with KIF) receive reduced
strength (0.40) and confidence (0.35), reflecting the lower
reliability of prose-level descriptions vs.\ formal axioms.
\paragraph{Measurement provenance note.}
All counts/percentages in this subsection come from repository-local scripts and
logs in the reported snapshot; they are engineering measurements, not Lean theorems.
Artifact pointers (snapshot): \texttt{Mettapedia/Languages/GF/SUMO/data/doc\_gf\_parse.jsonl},
\texttt{Mettapedia/Languages/GF/SUMO/data/doc\_gf\_parsed\_evidence\_atoms.jsonl},
\texttt{Mettapedia/Languages/GF/SUMO/data/doc\_evidence\_atoms.jsonl},
\texttt{Mettapedia/Languages/GF/SUMO/sumo\_doc\_gf\_parse.py},
\texttt{Mettapedia/Languages/GF/SUMO/sumo\_doc\_to\_evidence\_atoms.py}.

\paragraph{Five-lane evidence architecture.}
The full evidence model used in this pipeline slice draws from five operationally
independent lanes:
\begin{center}\small
\begin{tabular}{@{}lrlr@{}}
\toprule
\textbf{Lane} & \textbf{Weight} & \textbf{Independence group}
  & \textbf{Atoms} \\
\midrule
WM witnesses (\texttt{checkLang}) & 1.00 & \texttt{wm.*} & 6 \\
KIF rendered templates (\texttt{kifRendered}) & 0.50 & \texttt{kif.*} & 805 \\
GF runtime parsed (\texttt{engGFParsed}) & 0.60 & \texttt{gfparse.*} & 514 \\
Doc GF parsed (\texttt{docGFParsed}) & 0.50 & \texttt{docgf.*} & 60 \\
Doc English hints (\texttt{englishDoc}) & 0.30 & \texttt{doc.*} & 1{,}166 \\
\midrule
\textbf{Total} & & & \textbf{2{,}551} \\
\bottomrule
\end{tabular}
\end{center}

Independence groups prevent double-counting within and across
lanes: each group retains only the highest-weight atom.
Of the 26~repair decisions, 14~now have evidence atoms from
at least two independent lanes, enabling meaningful Beta
aggregation.  The \texttt{SumoNTT} regression suite remains
at 20~passed, 0~failed, 8~N/A (encoding choices and
meta-decisions not testable as reachability queries).
Artifact pointers (snapshot): \texttt{Mettapedia/Languages/GF/SUMO/data/pipeline\_audit.json},
\texttt{Mettapedia/Languages/GF/SUMO/data/kif\_rendered\_evidence\_atoms.jsonl},
\texttt{Mettapedia/Languages/GF/SUMO/data/gf\_parsed\_evidence\_atoms.jsonl},
\texttt{Mettapedia/Languages/GF/SUMO/data/doc\_gf\_parsed\_evidence\_atoms.jsonl},
\texttt{Mettapedia/Languages/GF/SUMO/data/doc\_evidence\_atoms.jsonl}.

%% ─────────────────────────────────────────────────────────────
\section{Comparison with Direct SUMO $\to$ Lean}
\label{sec:comparison}

We developed two direct approaches before arriving at the pipeline.

\subsection{sumo-lean4: KIF Checker and HOL Emitter}

In the reported snapshot, the \texttt{sumo-lean4} tool parses all 113 \KIF{} files and:
\begin{enumerate}[nosep]
  \item Checks arity consistency (relation usage vs.\ declaration).
  \item Emits ``HOL-native'' Lean code with curried predicates for
    fixed-arity relations and \texttt{List}-based application for
    variable-arity ones.
\end{enumerate}

After one pass in that snapshot, the checker reported ``no issues found'' across all
files (after applying the fixes in \S\ref{sec:sumo}).

\textbf{Strength}: Found genuine arity and argument bugs quickly.

\textbf{Limitation}: The emitted Lean code declares \SUMO{} axioms
as Lean \texttt{axiom}s---opaque to Lean's type system.
Once imported, bugs \emph{become trusted axioms in Lean's kernel}.
There is no behavioral or evidential checking.

\subsection{LeanSUMO: Type-Based Attributes}

The LeanSUMO project encoded \SUMO{} concepts as Lean structures with
attribute fields:
\begin{verbatim}
  structure Human extends Organism where
    sex : SexAttr
  structure Man extends Human where
    isMale : sex = SexAttr.Male
\end{verbatim}

This enabled elegant 4-line disjointness proofs (Man $\neq$ Woman)
and caught the Pleasure/Pain type confusion at compile time.

\textbf{Strength}: Lean's type system catches sort errors as
type mismatches.  Proofs are trivial.

\textbf{Limitation}: Each \SUMO{} class requires manual encoding
as a Lean structure.  With 20\,000+ classes, this does not scale.
No automation exists for the \KIF{} $\to$ structure translation.
And crucially, there is no behavioral semantics: the encoding
captures \emph{what things are} but not \emph{how they behave}.

\subsection{The Pipeline Approach}

The pipeline addresses both limitations:
\begin{itemize}[nosep]
  \item \textbf{Automation}: the \GF{} grammar is generated from
    \SUMO{}'s \texttt{domain}/\texttt{range}/\texttt{subclass}
    declarations, just as the HOL emitter already does.
    No manual per-class encoding.
  \item \textbf{Behavioral semantics}: \OSLF{} synthesis derives
    $\DiamondOp/\BoxOp$ types automatically from operational rules.
  \item \textbf{Evidence grounding}: the \WM{} layer quantifies how
    well-supported each assertion is.
  \item \textbf{Sort safety}: \GF{} categories provide the same sort
    discipline as LeanSUMO's typed attributes, but generated.
\end{itemize}

\begin{table}[ht]
\centering
\caption{Bug classes caught by each approach.}
\label{tab:comparison}
\small
\begin{tabular}{@{}lcccc@{}}
\toprule
\textbf{Bug class} & \textbf{sumo-lean4} & \textbf{LeanSUMO} & \textbf{Pipeline} \\
\midrule
Arity mismatch & \checkmark & \checkmark & \checkmark \\
Missing arguments & \checkmark & & \checkmark \\
Type/sort confusion & & \checkmark & \checkmark \\
Semantic contradiction & & & \checkmark \\
Behavioral inconsistency & & & \checkmark \\
Ungroundable assertion & & & \checkmark \\
Scope ambiguity & & & \checkmark \\
Compositional bug & & & \checkmark \\
\midrule
Repair disambiguation & & & \checkmark \\
\bottomrule
\end{tabular}
\end{table}

%% ─────────────────────────────────────────────────────────────
\section{Worked Example: The Classmate Arity Bug}
\label{sec:example}

We trace a real \SUMO{} bug through the pipeline.

\paragraph{The bug.}
\SUMO{} declares:
\begin{verbatim}
  (instance classmate IrreflexiveRelation)
  (domain classmate 1 Human)
  (domain classmate 2 Human)
  (domain classmate 3 Organization)
\end{verbatim}
But \texttt{IrreflexiveRelation} is defined in \texttt{Merge.kif} as
a subclass of \texttt{BinaryRelation}.  A ternary relation cannot be
an instance of a binary-only relation class.

\paragraph{Stage 1 (\SUMO{} $\to$ \GF{}).}
The \GF{} grammar generator creates:
\begin{verbatim}
  cat Human ; cat Organization ;
  fun classmate : Human -> Human -> Organization -> Prop ;
\end{verbatim}
and separately:
\begin{verbatim}
  fun IrreflexiveRelation_inst :
    (Entity -> Entity -> Prop) -> Prop ;
\end{verbatim}
The assertion \texttt{(instance classmate IrreflexiveRelation)} requires
passing \texttt{classmate} (a ternary function) to
\texttt{IrreflexiveRelation\_inst} (which expects a binary function).
\emph{The \GF{} type checker rejects this: type mismatch.}

\paragraph{Stage 2 (\GF{} $\to$ \OSLF{}).}
Had the bug not been caught at Stage~1, the \OSLF{} synthesis would
derive rewrite rules for \texttt{classmate} with three arguments.
The irreflexivity axiom
\texttt{(forall (?R) (=> (instance ?R IrreflexiveRelation) (not (?R ?X ?X))))}
would apply to the binary projection of \texttt{classmate}, producing
$\BoxOp(\neg\texttt{classmate}(a,a))$---but this is ill-formed because
\texttt{classmate} requires three arguments.
The $\BoxOp$ evaluation would flag an arity error in the behavioral
type.

\paragraph{Stage 3 (\OSLF{} $\to$ \WM{}).}
The evidence evaluator would assign $\ev{0}{0}$ (zero evidence) to
the irreflexivity claim for \texttt{classmate}, since the claim is
about a binary usage that never occurs in the ontology's operational
semantics.
This zero-evidence diagnostic would flag the assertion as
ungroundable.

\paragraph{Comparison.}
\begin{itemize}[nosep]
  \item \textbf{KIF checker}: catches the arity mismatch directly
    (our \texttt{sumo-lean4} found this).
  \item \textbf{ATP}: may or may not notice, depending on how the
    ternary/binary conflict interacts with other axioms.
  \item \textbf{OntoClean}: irrelevant (this is a syntactic/semantic
    bug, not a metaproperty issue).
  \item \textbf{Keet testing}: would generate test cases for
    irreflexivity, which might succeed or fail depending on the
    \ATP{}'s handling of the arity conflict.
  \item \textbf{Pipeline}: catches at Stage~1 (type error),
    Stage~2 ($\BoxOp$ ill-formedness), and Stage~3 (zero evidence).
    Three independent diagnostics from different layers.
\end{itemize}

%% ─────────────────────────────────────────────────────────────
\section{Related Work and Discussion}
\label{sec:related}

\paragraph{Faithful HOL embeddings.}
Benzm\"uller et al.~\cite{benzmuller-faithful-embeddings} embed
various logics (modal, deontic, conditional) into classical \HOL{},
enabling \ATP{} tools to reason about non-classical ontological
assertions.  Brown, Pease, and Urban~\cite{brown2022sumo-set-theory}
extended this by embedding the SUMO-K fragment into higher-order set
theory via Isabelle/HOL, providing formal semantics for \SUMO{}'s
higher-order features.
Our pipeline is complementary: where these embeddings
preserve the \emph{deductive} content of axioms, our pipeline extracts
their \emph{operational} content.
Combining both---deductive verification via faithful embedding plus
behavioral verification via OSLF---would provide broader end-to-end
coverage across deductive and behavioral checks.

\paragraph{Ontology repair.}
Troquard et al.~\cite{troquard2018weakening} introduced axiom weakening
as an alternative to deletion-based repair, using refinement operators
from inductive logic programming.  Baader and
Kriegel~\cite{baader2023pseudocontraction,baader2022fixedpremise}
established a formal correspondence between optimal repair and AGM
belief revision (pseudo-contraction), with practical algorithms for
$\mathcal{EL}$ TBoxes.  Alvez et al.~\cite{alvez2012adimensumo}
applied \ATP{}-based diagnosis specifically to \SUMO{}, re-engineering
88\% of the ontology.  Our repair loop (\S\ref{sec:repair}) extends
this line of work to full \FOL{} with behavioral and evidential
disambiguation signals that these approaches lack.

\paragraph{Autoformalization.}
The repair loop's structure---formalize, check, use error feedback to
correct, retry---mirrors the LLM autoformalization
loop~\cite{autoformalization-survey}, where a language model generates
formal statements and a proof assistant provides verification feedback.
The key difference is the richness of the feedback: our pipeline
provides type context, behavioral diagnostics, and evidence scores,
whereas autoformalization typically receives only type error messages.

\paragraph{Modular ontologies.}
Stuckenschmidt et al.~\cite{stuckenschmidt-modular-ontologies}
develop theory and techniques for decomposing large ontologies into
modules with well-defined interfaces.
Our \GF{}-based approach naturally supports modularity: each \SUMO{}
domain file becomes a \GF{} module, and the \OSLF{} synthesis respects
module boundaries.

\paragraph{Connection to Mettapedia.}
The referenced \GF{} $\to$ \OSLF{} $\to$ \WM{} pipeline slice is formalized in
the Mettapedia Lean~4 library and currently builds without
\texttt{sorry}/axiom placeholders in that slice.
The PLN evidence quantale, the Galois connection, the compositional
semantics, and the world-model bridge are represented there by theorem-level
endpoints.
\SUMO{}-\GF{} is a new concrete syntax for this existing,
verified infrastructure, and the Evidence Model
(\texttt{EvidenceModel.lean}) extends the repair log with structured
assurance scoring that connects directly to
\texttt{EvidenceQuantale}'s $(n^+, n^-)$ algebra.

\paragraph{Limitations.}
\begin{itemize}[nosep]
  \item The pipeline does not replace OntoClean-style adequacy
    judgments.  Philosophical metaproperties (rigidity, identity)
    require human expertise.
  \item The behavioral-correctness layer (\S\ref{def:behavioral})
    presupposes an operational reading of the ontology.  For purely
    classificatory fragments without process axioms, Stages~2--3
    add little value.
  \item The \SUMO{} $\to$ \GF{} translation (Stage~1) is not yet
    implemented as an automated tool; it is described as a
    well-motivated extension of the existing \texttt{sumo-lean4}
    HOL emitter.
  \item Row variables in \KIF{} require special handling (bounded
    unfolding or higher-order \GF{} categories).
  \item The formalization-driven repair loop (\S\ref{sec:repair}) is
    partially implemented: the detection pipeline, evidence model,
    and multi-agent review infrastructure are in place, with 26
    decisions logged and 20 passing formal verification.  The
    repair-selection heuristics and the top-down traversal strategy
    require further engineering effort and empirical validation on
    \SUMO{}'s full 20k axioms.  The \texttt{score\_repairs.py}
    tool to compute assurance numbers automatically per decision
    is planned but not yet built.
\end{itemize}
Artifact pointers (snapshot): \texttt{Mettapedia/Languages/GF/SUMO/RepairLog.lean},
\texttt{Mettapedia/Languages/GF/SUMO/SumoNTT.lean},
\texttt{Mettapedia/Languages/GF/SUMO/data/decision\_coverage.json},
\texttt{Mettapedia/Languages/GF/SUMO/data/review\_ledger\_2026-02-24.jsonl}.

%% ─────────────────────────────────────────────────────────────
\section{Conclusion}
\label{sec:conclusion}

Ontology correctness is not a single property but a hierarchy:
syntactic well-formedness, semantic consistency, taxonomic adequacy,
and behavioral correctness.
Existing methods address the first three layers but leave behavioral
adequacy---whether an ontology's axioms produce the intended
operational consequences---unexamined.

The pipeline \SUMO{} $\to$ \GF{} $\to$ \OSLF{} $\to$ \WM{} provides
the missing layers:
\begin{itemize}[nosep]
  \item \GF{}'s typed abstract syntax catches sort and arity errors
    at grammar level.
  \item \OSLF{}'s behavioral types ($\DiamondOp$, $\BoxOp$, native
    types) detect operational inconsistencies.
  \item The world-model's evidence-valued semantics flags ungroundable
    assertions, contradictory evidence flows, and scope ambiguities.
\end{itemize}

Beyond detection, the formalization-driven repair loop
(\S\ref{sec:repair}) provides a principled methodology for
\emph{correcting} ontology bugs: multi-level diagnostics disambiguate
among repair candidates, canary theorems prevent regression, and the
loop converges by construction.  The Evidence Model
(\S\ref{sec:evidence-model}) adds quantitative assurance:
Beta-distribution scoring with archetype-specific priors,
multi-agent review (Claude Code, Codex, GPT-5.2 Pro, human),
and PLN-compatible evidence atoms that integrate with the
proven evidence quantale.  No repair is applied until structured
evidence, assurance thresholds, and review quorum are satisfied.
This addresses the gap left by
SUMO-GF~\cite{enache2012sumogf}, whose type checker rejected 12\% of
axioms outright (with a further 24\% inexpressible), and which made
silent classification repairs with no audit trail or verification
methodology.

The \GF{} $\to$ \OSLF{} $\to$ \WM{} core is formalized in Lean~4 in the
referenced pipeline slice, with no \texttt{sorry}/axiom placeholders in that
slice.
Extending the pipeline to \SUMO{}---a large, openly available formal ontology---is
the natural next step: it inherits the benefits of both our
\texttt{sumo-lean4} arity checker and LeanSUMO's type-based
attributes, while adding behavioral and evidential correctness
layers that neither provides.  As a proof of concept, we encode the
FOET-relevant fragment of \SUMO{}-\GF{} and run it through the full
pipeline (Appendix~\ref{app:sumo-gf-experiment}), automatically
detecting known bugs including the Pain/Attribute type confusion.
The repair loop transforms the pipeline from a static analysis tool
into an interactive ontology engineering methodology.

%% ─────────────────────────────────────────────────────────────
\begin{thebibliography}{99}

\bibitem{pease2002sumo}
A.~Pease, I.~Niles, and J.~Li,
``The Suggested Upper Merged Ontology: A Large Ontology for the Semantic
Web and its Applications,''
in \textit{Working Notes of the AAAI-2002 Workshop on Ontologies and the
Semantic Web}, 2002.

\bibitem{pease-casc-prizes}
A.~Pease,
``The Annual SUMO Reasoning Prizes at CASC,''
\textit{CADE Workshop}, 2015.

\bibitem{pease-sutcliffe-2007}
A.~Pease and G.~Sutcliffe,
``First Order Reasoning on a Large Ontology,''
in \textit{Proc.\ ESARLT}, pp.~61--70, 2007.

\bibitem{pease-reasoning-generation}
A.~Pease,
``Reasoning and Language Generation in the SUMO Ontology,''
\textit{Technical Report}, 2019.

\bibitem{hayes-menzel-kif-2001}
P.~Hayes and C.~Menzel,
``A Semantics for the Knowledge Interchange Format,''
in \textit{Proc.\ IJCAI Workshop on the KIF Standard}, 2001.

\bibitem{kalyanpur2005}
A.~Kalyanpur, B.~Parsia, E.~Sirin, and J.~Hendler,
``Debugging Unsatisfiable Classes in OWL Ontologies,''
\textit{Journal of Web Semantics}, vol.~3, pp.~268--293, 2005.

\bibitem{keet2017whitebox}
C.~M.~Keet and A.~\L{}awrynowicz,
``Automatic White-Box Testing of First-Order Logic Ontologies,''
\textit{arXiv:1705.10219}, 2017.

\bibitem{keet2017blackbox}
C.~M.~Keet and A.~\L{}awrynowicz,
``Black-Box Testing of First-Order Logic Ontologies Using WordNet,''
\textit{arXiv:1705.10217}, 2017.

\bibitem{guarino-welty-ontoclean}
N.~Guarino and C.~A.~Welty,
``An Overview of OntoClean,''
in \textit{Handbook on Ontologies}, Springer, pp.~151--159, 2009.

\bibitem{gangemi-presutti-patterns}
A.~Gangemi and V.~Presutti,
``Ontology Design Patterns,''
in \textit{Handbook on Ontologies}, Springer, pp.~221--243, 2009.

\bibitem{ranta2004gf}
A.~Ranta,
``Grammatical Framework: A Type-Theoretical Grammar Formalism,''
\textit{Journal of Functional Programming}, vol.~14, no.~2, pp.~145--189,
2004.

\bibitem{ranta-typeful-ontologies}
A.~Ranta, K.~Angelov, and R.~Enache,
``Typeful Ontologies with Direct Multilingual Verbalization,''
in \textit{Proc.\ Controlled Natural Language}, 2012.

\bibitem{benzmuller-faithful-embeddings}
C.~Benzm\"uller,
``Faithful Logic Embeddings in HOL---Deep and Shallow,''
\textit{arXiv:2502.19311}, 2025.

\bibitem{stuckenschmidt-modular-ontologies}
H.~Stuckenschmidt, C.~Parent, and S.~Spaccapietra, Eds.,
\textit{Modular Ontologies: Concepts, Theories and Techniques for
Knowledge Modularization},
LNCS 5445, Springer, 2009.

\bibitem{goertzel2008pln}
B.~Goertzel, M.~Iklé, I.~F.~Goertzel, and A.~Heljakka,
\textit{Probabilistic Logic Networks: A Comprehensive Framework for
Uncertain Inference},
Springer, 2008.

\bibitem{bittner-formal-ontologies}
T.~Bittner,
``On the Computational Realization of Formal Ontologies,''
\textit{Technical Report}, University at Buffalo, 2004.

\bibitem{euzenat-shvaiko-matching}
J.~Euzenat and P.~Shvaiko,
\textit{Ontology Matching}, 2nd ed., Springer, 2013.

\bibitem{translating-sumo-hol}
A.~Pease, G.~Sutcliffe, N.~Siegel, and S.~Trac,
``Translating SUMO-K to Higher-Order Set Theory,''
2010.

\bibitem{troquard2018weakening}
N.~Troquard, R.~Confalonieri, P.~Galliani, R.~Pe\~naloza,
D.~Porello, and O.~Kutz,
``Repairing Ontologies via Axiom Weakening,''
in \textit{Proc.\ AAAI}, pp.~1981--1988, 2018.

\bibitem{baader2023pseudocontraction}
F.~Baader and F.~Kriegel,
``Optimal Repairs in Ontology Engineering as Pseudo-Contractions
in Belief Change,''
in \textit{Proc.\ ACM SAC}, pp.~846--855, 2023.

\bibitem{baader2022fixedpremise}
F.~Baader and F.~Kriegel,
``Optimal Fixed-Premise Repairs of $\mathcal{EL}$ TBoxes,''
in \textit{Proc.\ KI}, LNCS 13404, pp.~103--117, Springer, 2022.

\bibitem{alvez2012adimensumo}
J.~Alvez, P.~Lucio, and G.~Rigau,
``Adimen-SUMO: Reengineering an Ontology for First-Order Reasoning,''
\textit{International Journal on Semantic Web and Information Systems},
vol.~8, no.~4, pp.~80--116, 2012.

\bibitem{enache2012sumogf}
R.~Enache and K.~Angelov,
``Typeful Ontologies with Direct Multilingual Verbalization,''
in \textit{Proc.\ Controlled Natural Language}, LNCS 7427,
pp.~52--67, Springer, 2012.

\bibitem{brown2022sumo-set-theory}
C.~E.~Brown, A.~Pease, and J.~Urban,
``Translating SUMO-K to Higher-Order Set Theory,''
in \textit{Proc.\ FroCoS}, LNCS 14279, pp.~253--274, Springer, 2023.

\bibitem{horridge2008laconic}
M.~Horridge, B.~Parsia, and U.~Sattler,
``Laconic and Precise Justifications in OWL,''
in \textit{Proc.\ ISWC}, LNCS 5318, pp.~323--338, Springer, 2008.

\bibitem{schlobach2003diagnosis}
S.~Schlobach and R.~Cornet,
``Non-Standard Reasoning Services for the Debugging of Description
Logic Terminologies,''
in \textit{Proc.\ IJCAI}, pp.~355--362, 2003.

\bibitem{autoformalization-survey}
Y.~Wu, Z.~Jiang, et al.,
``Autoformalization in the Era of Large Language Models: A Survey,''
\textit{arXiv:2505.23486}, 2025.

\bibitem{foet2022}
B.~Goertzel and Z.~Kothare,
``The Formal Ethics Seed Ontology (FOET),''
\textit{Garden of Minds}, 2022.
\url{https://gardenofminds.art/ethics/formal-ethics-seed-ontology/}

\bibitem{benzmuller-gewirth}
C.~Benzm\"uller, A.~Fuenmayor, and D.~Fuenmayor,
``Formalisation and Evaluation of Alan Gewirth's Proof for the
Principle of Generic Consistency in Isabelle/HOL,''
\textit{Archive of Formal Proofs}, 2018.

\end{thebibliography}

%% ─────────────────────────────────────────────────────────────
\appendix

\section{Case Study: SUMO Dependencies of the Formal Ethics Seed Ontology}
\label{app:foet}

To make the repair loop concrete, we identify the \emph{exact} \SUMO{}
fragment required by a real downstream application: the Formal Ethics
Seed Ontology (FOET)~\cite{foet2022}, a 6\,853-line \KIF{}
formalization of four ethical paradigms (deontological, value-judgment,
utilitarian, virtue ethics) with inter-paradigm translations and a
Lean~4 formalization including a fully proven Gewirth PGC
argument~\cite{benzmuller-gewirth}.
This appendix traces FOET's \SUMO{} dependencies upward through the
class hierarchy, producing a minimal subontology that must be fully
debugged and repaired before the pipeline can verify FOET end-to-end.

\subsection{FOET's SUMO Concept Inventory}

The FOET ontology introduces domain-specific classes (left column) that
are positioned within \SUMO{}'s existing hierarchy (right column).
Table~\ref{tab:foet-concepts} lists the key concept groups.

\begin{table}[ht]
\centering
\caption{FOET concepts and their \SUMO{} superclass anchors.}
\label{tab:foet-concepts}
\small
\begin{tabular}{@{}p{4.2cm}p{4cm}p{4cm}@{}}
\toprule
\textbf{FOET concept} & \textbf{Direct superclass} & \textbf{SUMO anchor} \\
\midrule
\multicolumn{3}{@{}l}{\emph{Ethical attributes}} \\
\texttt{MoralValueAttribute} & \texttt{MoralAttribute} & \texttt{NormativeAttribute} \\
\texttt{DeonticAttribute} & \texttt{ObjectiveNorm}, \texttt{MoralAttribute} & \texttt{NormativeAttribute} \\
\texttt{VirtueAttribute} & \texttt{MoralVirtueAttribute} & \texttt{PsychologicalAttribute} \\
\texttt{ViceAttribute} & \texttt{MoralVirtueAttribute} & \texttt{PsychologicalAttribute} \\
\midrule
\multicolumn{3}{@{}l}{\emph{Agents and processes}} \\
\texttt{VirtuousAgent} & \texttt{AutonomousAgent} & \texttt{Object} \\
\texttt{AutonomousAgentProcess} & \texttt{Process} & \texttt{Physical} \\
\texttt{VirtuousAct} & \texttt{AutonomousAgentProcess} & \texttt{Process} \\
\texttt{ChoicePoint} & \texttt{Set}, \texttt{NonNullSet} & \texttt{SetOrClass} \\
\texttt{Situation} & (direct) & \texttt{Physical} \\
\midrule
\multicolumn{3}{@{}l}{\emph{Propositions and theories}} \\
\texttt{ActionFormula} & \texttt{Formula} & \texttt{Sentence} \\
\texttt{Conjecture} & \texttt{Sentence} & \texttt{LinguisticExpression} \\
\texttt{EthicalTheory} & (set of \texttt{EthicalSentence}) & \texttt{Proposition} \\
\texttt{Argument} & \texttt{Proposition} & \texttt{Abstract} \\
\midrule
\multicolumn{3}{@{}l}{\emph{Values (Schwartz--Haidt)}} \\
\texttt{Value} & (direct) & \texttt{Abstract} \\
10 Schwartz values & \texttt{Value} & \texttt{Abstract} \\
\midrule
\multicolumn{3}{@{}l}{\emph{Utility}} \\
\texttt{UtilityFormulaFn} & \texttt{UnaryFunction} & \texttt{Function} \\
\texttt{Consequentialism} & \texttt{Ethics} & domain module \\
\bottomrule
\end{tabular}
\end{table}

\subsection{SUMO Relations Used by FOET}

FOET introduces 30+ relations.  The most critical (those appearing in
axioms with inter-paradigm translations) are:

\begin{center}
\small
\begin{tabular}{@{}lll@{}}
\toprule
\textbf{Relation} & \textbf{SUMO type} & \textbf{Key domains} \\
\midrule
\texttt{hasAttribute} & \texttt{BinaryPredicate} & Agent $\times$ Attribute \\
\texttt{desires} & \texttt{PropositionalAttitude} & Agent $\times$ Formula \\
\texttt{holdsValue} & \texttt{BinaryPredicate} & Agent $\times$ Value \\
\texttt{holdsObligation} & \texttt{BinaryPredicate} & Formula $\times$ Agent \\
\texttt{realizes\-Formula} & \texttt{BinaryPredicate} & Process $\times$ Formula \\
\texttt{capableInSituation} & \texttt{QuaternaryPredicate} & ProcessClass $\times$ CaseRole $\times$ Object $\times$ Situation \\
\texttt{situationFn} & \texttt{UnaryFunction} & Process $\to$ Situation \\
\texttt{influences} & \texttt{BinaryPredicate} & (general) \\
\texttt{conjectures} & \texttt{PropositionalAttitude} & Agent $\times$ Conjecture \\
\bottomrule
\end{tabular}
\end{center}

\subsection{Upward Dependency Graph}

Figure~\ref{fig:foet-deps} shows the \SUMO{} class hierarchy from
\texttt{Entity} down to FOET's leaf concepts.  \emph{Every node on a
path from \texttt{Entity} to a FOET concept must be correctly
formalized} for the pipeline to verify FOET's axioms.

\begin{figure}[ht]
\centering
\small
\begin{verbatim}
Entity
+-- Physical
|   +-- Object
|   |   +-- SelfConnectedObject
|   |   |   `-- CorpuscularObject
|   |   |       `-- ContentBearingObject
|   |   `-- AutonomousAgent <- VirtuousAgent, ViciousAgent
|   |       `-- SentientAgent
|   |           `-- CognitiveAgent
|   +-- Process <- AutonomousAgentProcess
|   |   +-- VirtuousAct, ViciousAct
|   |   +-- BodyMotion, Vocalizing
|   |   `-- Change
|   |   `-- LinguisticCommunication
|   |       `-- Conjecturing
|   +-- ContentBearingPhysical
|   |   `-- LinguisticExpression
|   |       `-- Sentence
|   |           `-- Formula <- ActionFormula, Conjecture
|   `-- Situation [FOET-defined]
|
`-- Abstract
    +-- Attribute
    |   +-- InternalAttribute
    |   |   `-- BiologicalAttribute
    |   |       `-- PsychologicalAttribute
    |   |           +-- VirtueAttribute [FOET]
    |   |           `-- ViceAttribute [FOET]
    |   `-- RelationalAttribute
    |       `-- NormativeAttribute
    |           +-- ObjectiveNorm
    |           |   `-- DeonticAttribute [FOET]
    |           `-- MoralAttribute [FOET]
    |               +-- MoralValueAttribute [FOET]
    |               `-- MoralVirtueAttribute [FOET]
    +-- SetOrClass
    |   +-- Class
    |   `-- Set <- ChoicePoint
    +-- Relation
    |   +-- Predicate (BinaryPredicate, TernaryPredicate, ...)
    |   +-- Function (UnaryFunction, BinaryFunction)
    |   |   `-- UtilityFormulaFn [FOET]
    |   `-- InheritableRelation
    |       `-- PropositionalAttitude (desires, conjectures)
    +-- Proposition
    |   `-- Argument <- ValidDeductiveArgument
    |       `-- ConsequentialistArgument [FOET]
    +-- List <- SentenceList [FOET]
    `-- Value [FOET] <- 10 Schwartz values
\end{verbatim}
\caption{SUMO class hierarchy from \texttt{Entity} to FOET leaf
  concepts.  Nodes marked [FOET] are introduced by the ethics ontology;
  all others are \SUMO{} core (Merge.kif).  Every path from root to
  leaf must be fully debugged.}
\label{fig:foet-deps}
\end{figure}

\subsection{Repair Scope Estimate}

From the dependency graph, we identify three ``repair strata'' ordered
by priority:

\begin{enumerate}
  \item \textbf{Stratum 0 (Upper ontology core)}: \texttt{Entity},
    \texttt{Physical}, \texttt{Abstract}, \texttt{Object},
    \texttt{Process}, \texttt{Attribute}, \texttt{SetOrClass},
    \texttt{Relation}, \texttt{Proposition}.
    These 9 classes and their inter-relations (subclass, instance,
    domain, range) form the skeleton.  Approximately 200 axioms in
    Merge.kif.  Known bugs: minimal (heavily tested by
    Adimen-SUMO~\cite{alvez2012adimensumo}).

  \item \textbf{Stratum 1 (Mid-level anchors)}: \texttt{Agent} chain
    (\texttt{AutonomousAgent} $\to$ \texttt{SentientAgent} $\to$
    \texttt{CognitiveAgent}), \texttt{NormativeAttribute} $\to$
    \texttt{ObjectiveNorm}, \texttt{PsychologicalAttribute},
    \texttt{Formula}/\texttt{Sentence}, \texttt{Predicate}/\texttt{Function},
    \texttt{PropositionalAttitude}.
    Approximately 500 axioms.  Known bugs: type confusion in the
    \texttt{Attribute} hierarchy (Pain/Pleasure,
    \S\ref{sec:sumo}); potential arity issues in
    \texttt{PropositionalAttitude} predicates.

  \item \textbf{Stratum 2 (FOET-specific)}: \texttt{MoralAttribute},
    \texttt{DeonticAttribute}, \texttt{VirtueAttribute},
    \texttt{AutonomousAgentProcess}, \texttt{ChoicePoint},
    \texttt{Situation}, \texttt{Value}, \texttt{UtilityFormulaFn},
    and all 30+ FOET relations.
    Approximately 300 axioms (from the FOET \KIF{} file).
    Known bugs: the FOET \KIF{} itself documented several
    \SUMO{} compatibility issues in
    \texttt{ResolvedBugsInSUMO.kif}.
\end{enumerate}
Stratum-count provenance (snapshot): \texttt{Mettapedia/Languages/GF/SUMO/stratum\_scan.py},
\texttt{Mettapedia/Languages/GF/SUMO/stratum0\_scan\_report.txt},
\texttt{Mettapedia/Languages/GF/SUMO/stratum1\_scan\_with\_foet\_report.txt},
\texttt{Mettapedia/Languages/GF/SUMO/stratum2\_scan\_with\_foet\_report.txt}.

Total: approximately \textbf{1\,000 axioms} spanning 3 strata---a
tractable target for the repair loop.  This is 5\% of \SUMO{}'s full
20k axioms, but it covers the upward-closure slice targeted for the
ethics-application fragment in this draft.

\subsection{Why FOET Is a Good First Target}

\begin{enumerate}[nosep]
  \item \textbf{Lean formalization exists}: the FOET Lean~4 project
    (27 modules, 4\,607 lines) already type-checks with a minimal
    \SUMO{} kernel (\texttt{SumoCore.lean}: \texttt{Class},
    \texttt{instanceOf}, \texttt{subclass}).  The repair loop would
    upgrade this kernel to the full 1\,000-axiom fragment.
  \item \textbf{Gewirth PGC is proven}: the culminating theorem
    (Principle of Generic Consistency) is proven in both Isabelle/HOL
    and Lean~4 without \texttt{sorry}.  Canary theorems for the
    repair loop can be drawn directly from this proof.
  \item \textbf{Inter-paradigm translations}: the proven loop
    (Value $\to$ Deontic $\to$ Utilitarian $\to$ Virtue $\to$ Value
    $=$ id) provides a rich set of canary theorems spanning multiple
    \SUMO{} concept groups.
  \item \textbf{Practical relevance}: formal ethics for AI alignment
    requires certified reasoning over normative concepts---exactly the
    use case where ontology correctness matters most.
\end{enumerate}

%% ─────────────────────────────────────────────────────────────
\section{Experiment: SUMO-GF Through the Pipeline}
\label{app:sumo-gf-experiment}

We encode the FOET-relevant fragment of the \SUMO{}-\GF{} abstract
syntax~\cite{enache2012sumogf} in Lean~4 and run it through the full
\GF{} $\to$ \OSLF{} $\to$ \WM{} pipeline, using exactly the same
infrastructure as the English and Czech grammars.

\subsection{Encoding Strategy: Shallow Flattening}

The \SUMO{}-\GF{} abstract syntax uses GF's dependent types:
\begin{verbatim}
  cat Class; El Class; Ind Class; Var Class;
      SubClass (c1,c2 : Class); Inherits Class Class;
      Formula; Stmt;
\end{verbatim}
Our pipeline uses flat string-valued sorts.  We apply a
\emph{shallow flattening}: for each of the ${\sim}50$ FOET-relevant
\SUMO{} classes $C$, we generate concrete sort names:
\[
  \texttt{El}~C \;\mapsto\; \texttt{"El\_}C\texttt{"},\qquad
  \texttt{Ind}~C \;\mapsto\; \texttt{"Ind\_}C\texttt{"},\qquad
  \texttt{Var}~C \;\mapsto\; \texttt{"Var\_}C\texttt{"}
\]
This loses GF's parametric polymorphism but preserves the
sort-distinction structure that makes bugs visible, and it is
compatible with the existing \texttt{gfGrammarLanguageDef}
infrastructure without modification.

The 50 classes are drawn from the FOET dependency graph
(Appendix~\ref{app:foet}), organized into three strata:
\begin{itemize}[nosep]
  \item \textbf{Stratum~0} (9 classes): \texttt{Entity}, \texttt{Physical},
    \texttt{Abstract}, \texttt{Object}, \texttt{Process},
    \texttt{Attribute}, \texttt{SetOrClass}, \texttt{Relation},
    \texttt{Proposition}
  \item \textbf{Stratum~1} (18 classes): the Agent chain, normative
    attributes, psychological attributes, linguistic expressions
  \item \textbf{Stratum~2} (23 classes): FOET-specific concepts
    (virtues, deontic attributes, choice points, utility functions)
\end{itemize}

\subsection{Pipeline Statistics}

Table~\ref{tab:sumo-gf-stats} compares the \SUMO{}-\GF{} encoding
with the existing RGL (English/Czech) encoding.

\begin{table}[h]
\centering
\begin{tabular}{lrr}
\toprule
& \textbf{RGL (Eng/Cz)} & \textbf{SUMO-GF} \\
\midrule
Categories (sorts) & 25 & ${\sim}200$ \\
Grammar rules & 170 & ${\sim}350$ \\
Rewrite rules & 14 & ${\sim}115$ \\
Equations & 2 & 50 \\
\midrule
Galois connection & proven & proven \\
Native type theory & yes & yes \\
\bottomrule
\end{tabular}
\caption{Pipeline statistics: RGL vs.\ \SUMO{}-\GF{} encoding.
  The \SUMO{}-\GF{} encoding is larger because each class generates
  separate sort, coercion, and quantifier entries.}
\label{tab:sumo-gf-stats}
\end{table}

The larger number of sorts and rules reflects \SUMO{}'s richer type
structure: where the RGL has a single \texttt{NP} sort, \SUMO{}-\GF{}
has a separate \texttt{El\_$C$} for each class $C$, with explicit
coercion functions between them.

\subsection{Rewrite Dynamics}

We define four classes of rewrite rules, mirroring the GF \texttt{data}/\texttt{def} structure:

\begin{enumerate}
  \item \textbf{Reflexive coercion elimination} (${\sim}100$ rules):
    $\texttt{el\_}C\texttt{\_}C(x) \;\Red\; x$ for each class $C$.
    In GF, \texttt{el(C, C, inhz(C), x) = x}.

  \item \textbf{Variable coercion elimination} (${\sim}50$ rules):
    $\texttt{var\_}C\texttt{\_}C(x) \;\Red\; x$.

  \item \textbf{Transitive chain collapse} (11 rules):
    $\texttt{el\_}B\texttt{\_}C(\texttt{el\_}A\texttt{\_}B(x))
      \;\Red\; \texttt{el\_}A\texttt{\_}C(x)$
    for selected chains $A <: B <: C$.

  \item \textbf{Double negation}: $\lnot\lnot\varphi \;\Red\; \varphi$.
\end{enumerate}

All four classes produce verified output from \texttt{\#eval!} in
Lean~4.  The reflexive and variable coercions are particularly
important: they ensure that $\DiamondOp$ is \emph{non-vacuous} for
identity coercions, meaning the behavioral type system has actual
content rather than being trivially satisfied.

\subsection{Bug Detection Results}

\paragraph{Pain/Attribute type confusion (detected and analyzed).}
In \SUMO{}'s \texttt{Mid-level-ontology.kif}, \texttt{Pain} is
declared \texttt{(subclass Pain PathologicProcess)}---a
\texttt{Process}.  Meanwhile, \texttt{contraryAttribute} expects
arguments of type \texttt{El\_Attribute}.  Our pipeline faithfully
follows \SUMO{}'s taxonomy:
\begin{itemize}[nosep]
  \item \texttt{Pain} : \texttt{Ind\_Process}
  \item \texttt{Pleasure} : \texttt{Ind\_Attribute}
  \item \texttt{contraryAttribute} : \texttt{El\_Attribute} $\to$
    \texttt{El\_Attribute} $\to$ \texttt{SumoFormula}
\end{itemize}
The pipeline detects that \texttt{Pain} has coercions only to
\texttt{El\_Process} $\to$ \texttt{El\_Physical} $\to$
\texttt{El\_Entity}, with \emph{no path} to \texttt{El\_Attribute}.

\paragraph{Comparison with the original \SUMO{}-\GF{}.}
Examination of \texttt{MidLevelOntology.gf} from the GF~3.3
distribution reveals that Enache silently reclassified Pain:
\begin{verbatim}
  fun Pain : Ind EmotionalState ;   -- MidLevelOntology.gf:4402
\end{verbatim}
where \texttt{EmotionalState} $\sqsubseteq$ \texttt{StateOfMind}
$\sqsubseteq$ \texttt{PsychologicalAttribute} $\sqsubseteq \cdots
\sqsubseteq$ \texttt{Attribute}.  This makes
\texttt{contraryAttribute(Pleasure, Pain)} type-correct by moving
\texttt{Pain} from the Process branch to the Attribute branch---an
\emph{undocumented repair}.

This repair is defensible: in ordinary language, pain is a
\emph{sensation} (Attribute) that is \emph{caused by} a pathological
process, not the process itself.  But the axiom usage census tells
a different story: 5 axioms use \texttt{Pain} in Process-like
contexts (\texttt{instance}, \texttt{located}, \texttt{WhenFn})
versus only 1 in an Attribute-like context
(\texttt{contraryAttribute}).  Raw evidence-counting would automate
to the \emph{wrong} classification.

The deeper judgment requires \emph{causal analysis} (Pain is the
effect, not the cause) and \emph{similarity-class evidence}
(\texttt{Pain} is used like \texttt{Pleasure}, which is
\texttt{EmotionalState}, so Pain should be too).  Our repair log
(\S\ref{sec:evidence-repair}) records this decision with evidence,
confidence, and canary theorems.

\paragraph{Behavioral dimension.}
Even if a permissive coercion \texttt{el\_Process\_Attribute} were
added, $\DiamondOp$ would test whether it has non-trivial rewrite
behavior.  A vacuous coercion (no reduction rules) would be flagged
as behaviorally empty---a diagnostic that the original \SUMO{}-\GF{}
cannot provide.

\paragraph{Classmate arity mismatch (explained).}
\SUMO{} declares \texttt{classmate} as a \texttt{TernaryPredicate}
(arity 3: student, student, class) but also declares it an
\texttt{IrreflexiveRelation}, whose axiom schema
$\forall x.\;\lnot R(x,x)$ assumes arity~2.  In the \SUMO{}-\GF{}
encoding, \texttt{classmate} has type:
\[
  \texttt{El\_Agent} \to \texttt{El\_Agent} \to \texttt{El\_Class}
    \to \texttt{SumoFormula}
\]
The binary irreflexivity pattern cannot match this ternary signature.
The repair (from \S\ref{sec:example}) is to replace the binary
axiom with the properly quantified:
$\forall a,c.\;\lnot\texttt{classmate}(a, a, c)$.

\subsection{Comparison with the Original \SUMO{}-\GF{}}
\label{sec:evidence-repair}

We obtained the original \SUMO{}-\GF{} source files from the GF~3.3
distribution (30\,483 lines across 18 \texttt{.gf} files, including
1\,722 \texttt{fun} declarations in \texttt{Merge.gf} and 2\,827 in
\texttt{MidLevelOntology.gf}; measured in our local copy of that distribution).
Comparing the original with our
replication reveals both \emph{silent repairs} and
\emph{design choices} in the original encoding that were not
discussed in the published work.

\paragraph{Silent repairs.}
The original \SUMO{}-\GF{} contains at least one undocumented
reclassification (\texttt{Pain}: \texttt{PathologicProcess} $\to$
\texttt{EmotionalState}).  Our pipeline, by faithfully following
\SUMO{}'s own declarations, \emph{detects} the conflict rather than
silently resolving it.  The repair loop (\S\ref{sec:repair}) then
provides a principled framework for making and documenting such
decisions.

\paragraph{Variable-arity relations.}
The original encoding correctly types \texttt{contraryAttribute}
as \texttt{[El Attribute] -> Formula} (a GF list type), respecting
\SUMO{}'s declaration as a \texttt{VariableArityRelation}.  Our
initial encoding simplified this to a binary signature---a design
choice that should be corrected.  This repair is fully automatable:
any relation declared \texttt{(instance R VariableArityRelation)}
should receive a list-type signature (a mechanically checkable transformation
in this encoding pipeline).

\paragraph{Unresolved inheritance proofs.}
The axiom files (\texttt{.gft}) contain 592 lines with unresolved
\texttt{?} metavariables (out of 1\,184 axiom lines)---these
represent coercion proofs that GF's type inference could not
automatically resolve.  The pattern
\texttt{var~?~?~?~VARIABLE} indicates that the class, target class,
and inheritance proof for a variable coercion are all unknown.
In our pipeline, these would be flagged as incomplete behavioral
types requiring either explicit coercion paths or evidence-based
inference.

\paragraph{Enache's actual statistics.}
Enache's thesis reports that 64\% of quantified axioms passed
type-checking~\cite{enache2012sumogf}.  The 36\% failure breaks down as:
\begin{itemize}[nosep]
  \item 12\% genuinely rejected by the type checker
  \item 12\% inexpressible (variable types from usage context)
  \item 11\% using untranslated predicates (\texttt{range},
    \texttt{domain})
  \item 1\% other
\end{itemize}
For the higher-order axioms (involving meta-properties like
\texttt{AsymmetricRelation}), only 43\% translated successfully,
largely because \texttt{AsymmetricRelation} in GF assumes same-type
arguments while \SUMO{} applies it to cross-type relations like
\texttt{frequency : Process $\to$ TimeDuration}~\cite{enache2012sumogf}.

\paragraph{Evidence-based classification automation.}
\label{par:evidence-classify}
Our analysis suggests a three-tier automation strategy for
classification judgments:

\begin{enumerate}[nosep]
  \item \textbf{Axiom usage census}: Count how a concept is used
    across all \SUMO{} axioms.  Process-like patterns:
    \texttt{instance}, \texttt{located}, \texttt{WhenFn},
    \texttt{agent}, \texttt{patient}.  Attribute-like patterns:
    \texttt{attribute}, \texttt{contraryAttribute},
    \texttt{property}.  This is fully automatable.

  \item \textbf{Similarity-class evidence}: If concept $X$ is used
    in the same axiom patterns as concept $Y$, and $Y$ has a known
    classification, then $X$ should likely share that classification.
    For Pain: \texttt{Pain} appears alongside \texttt{Pleasure} in
    \texttt{contraryAttribute}, and \texttt{Pleasure} is
    \texttt{EmotionalState}, so Pain should be too.  This is
    automatable via co-occurrence analysis.

  \item \textbf{Causal/literature analysis}: When census and
    similarity evidence conflict (as for Pain: 5 Process-like uses
    vs.\ 1 Attribute-like use, but similarity says Attribute),
    deeper judgment is needed.  An LLM or commonsense KB can
    distinguish ``$X$ is used as $Y$'' from ``$X$ \emph{is} $Y$''
    by identifying causal structure (Pain is \emph{caused by} a
    process but \emph{is} a sensation).
\end{enumerate}

Each classification decision is recorded in a structured
\emph{repair log} (implemented as a Lean~4 structure with fields
for original classification, evidence items, confidence, automation
method, and canary theorems).  This log serves as an audit trail
that the original \SUMO{}-\GF{} lacks.

\paragraph{What the pipeline adds.}
Beyond reproducing Enache's sort-level detections, our pipeline
provides:
\begin{enumerate}[nosep]
  \item \textbf{Rewrite dynamics}: The $\DiamondOp/\BoxOp$ operators
    test whether coercions and axioms have non-trivial operational
    content, not just type-correctness.
  \item \textbf{Proven Galois connection}: The adjunction
    $\DiamondOp \dashv \BoxOp$ is proven automatically for the
    \SUMO{}-\GF{} language definition (via \texttt{langGalois}).
  \item \textbf{Repair audit trail}: Every classification decision
    is logged with evidence, confidence, and canary theorems---making
    the encoding \emph{reproducible and auditable}, unlike the
    original's silent repairs.
  \item \textbf{Cross-grammar invariance}: The same logical form
    can be expressed in both the RGL (English) grammar and
    \SUMO{}-\GF{}.  Agreement in \OSLF{} type assignment provides a
    grammar-invariant correctness check.
\end{enumerate}

\subsection{Implementation}

The encoding is implemented in five Lean~4 files:
\begin{itemize}[nosep]
  \item \texttt{SumoAbstract.lean}: 57 class definitions, 54 subclass
    edges, ${\sim}350$ function signatures (logical connectives,
    quantifiers per class, coercions per edge, \SUMO{} relations)
  \item \texttt{SumoOSLFBridge.lean}: Rewrites, the
    \texttt{LanguageDef}, OSLF type system, Galois connection proof,
    and executable pipeline tests
  \item \texttt{SumoNTT.lean}: Class hierarchy as GSLT, NTT
    extraction, comprehensive audit (26 decisions, 20 testable)
  \item \texttt{EvidenceModel.lean}: Repair archetypes, evidence
    atoms, review voting, assurance scoring (282~lines)
  \item \texttt{RepairLog.lean}: 26 repair decisions with evidence,
    WM witnesses, and review state tracking
\end{itemize}
Count provenance (snapshot): generated from these modules plus
\texttt{Mettapedia/Languages/GF/SUMO/audit\_repair\_pipeline.py} and
\texttt{Mettapedia/Languages/GF/SUMO/data/pipeline\_audit.json}.
In the reported snapshot, the listed pipeline files build clean; this claim is
scoped to those files rather than the entire repository.  The Galois connection
$\DiamondOp \dashv \BoxOp$ for \SUMO{}-\GF{} is proven automatically
by the same \texttt{langGalois} lemma used for English and Czech.

\subsection{Class Hierarchy as GSLT: NTT Extraction}
\label{sec:ntt-extraction}

A key discovery is that the \SUMO{} class hierarchy \emph{is itself}
a graded semantic lattice theory (GSLT).  The subclass relation
\texttt{Human} $\sqsubseteq$ \texttt{CognitiveAgent} $\sqsubseteq$
\texttt{SentientAgent} $\sqsubseteq \cdots \sqsubseteq$
\texttt{Entity} is a rewrite system: each coercion
$\texttt{el\_}c_1\texttt{\_}c_2(x) \;\Red\; x$ models subsumption as
reduction.  With the lattice operations \texttt{both} (meet) and
\texttt{either} (join) from GF's \texttt{Basic.gf}, this forms a GSLT.

We construct a fresh \texttt{LanguageDef} (\texttt{sumoHierarchyLangDef})
where:
\begin{itemize}[nosep]
  \item \textbf{Sorts} = \SUMO{} class names (33 for Layer~1)
  \item \textbf{Terms} = lift constructors + coercion rules
    (from the transitive closure: 153 edges from 34 direct edges)
  \item \textbf{Rewrites} = subclass-as-reduction
    ($\texttt{lift\_}c_1(x) \;\Red\; \texttt{lift\_}c_2(x)$
    for each $c_1 \sqsubseteq c_2$)
\end{itemize}
Passing this through \texttt{langOSLF} automatically produces
$\DiamondOp \dashv \BoxOp$ and NTT\@.  The Galois connection is proven
by the same \texttt{langGalois} lemma.

\paragraph{Behavioral types.}
The \emph{behavioral type} of a concept $X$ is the set of sorts
reachable via subclass coercion:
\[
  \mathit{BT}(X) = \{S \mid \DiamondOp(\mathit{is\_}S)(X)\}
\]
For example:
\begin{align*}
  \mathit{BT}(\texttt{CognitiveAgent}) &= \{\texttt{CognitiveAgent},
    \texttt{SentientAgent}, \texttt{AutonomousAgent}, \texttt{Object},
    \texttt{Physical}, \texttt{Entity}\} \\
  \mathit{BT}(\texttt{Process}) &= \{\texttt{Process},
    \texttt{Physical}, \texttt{Entity}\} \\
  \mathit{BT}(\texttt{Attribute}) &= \{\texttt{Attribute},
    \texttt{Abstract}, \texttt{Entity}\}
\end{align*}

\paragraph{Automated gap detection.}
A \emph{type gap} occurs when a concept $X$ is used in an axiom
context demanding sort $S$, but $S \notin \mathit{BT}(X)$.
The NTT extraction on Layer~1 (strata 0--1) automatically
detects four gaps:

\begin{center}
\small
\begin{tabular}{@{}llll@{}}
\toprule
\textbf{Concept} & \textbf{Relation} & \textbf{Demanded sort}
  & \textbf{Fix} \\
\midrule
\texttt{Pain} & \texttt{contraryAttribute} & \texttt{Attribute}
  & split/reclassify \\
\texttt{VirtuousAgent} & \texttt{attribute} & \texttt{Attribute}
  & FOET modeling error \\
\texttt{(our sig)} & \texttt{capableInSituation} arg~2 & \texttt{CaseRole}
  & encoding fixed \\
\texttt{(our sig)} & \texttt{capableInSituation} arg~3 & \texttt{Object}
  & encoding fixed \\
\bottomrule
\end{tabular}
\end{center}

Gaps 3--4 were encoding errors in our replication (wrong KIF domain
readings); these are fully automatable and have been corrected.
Gap~2 is a genuine modeling error in FOET: \texttt{VirtuousAgent} is
a \emph{class} of agents ($\sqsubseteq$ \texttt{AutonomousAgent}),
not an \texttt{Attribute} value, yet FOET uses
\texttt{(attribute ?AGENT VirtuousAgent)} at line~4733 of the seed
ontology.  The correct SUMO pattern is to create an individual
attribute (cf.\ \texttt{(instance Pietas VirtueAttribute)} at the
same file's line~4739, correctly used as
\texttt{(attribute ?AGENT Pietas)}).
Gap~1 (Pain) requires causal/literature evidence as discussed above.

\paragraph{Variable-name type inference.}
\SUMO{} KIF variable names carry implicit type information:
\texttt{?AGENT} implies \texttt{AutonomousAgent},
\texttt{?PROC} implies \texttt{Process},
\texttt{?ATTR} implies \texttt{Attribute}, and so on.
While imprecise, these conventions enable \emph{type-directed
hole-filling}: given a type gap where concept $X$ fills a slot
demanding sort $S$, we compute the \emph{inverse behavioral type}
$\mathit{BT}^{-1}(S) = \{C \mid S \in \mathit{BT}(C)\}$---the
set of all concepts that could type-correctly fill the slot.
For example, the 7 inhabitants of \texttt{Attribute} in Layer~1
are: \texttt{Attribute}, \texttt{InternalAttribute},
\texttt{BiologicalAttribute}, \texttt{PsychologicalAttribute},
\texttt{RelationalAttribute}, \texttt{NormativeAttribute},
\texttt{ObjectiveNorm}.  \texttt{VirtuousAgent} is not among them,
confirming the type gap.

Combined with pattern matching against analogous correct usage
in the same file (e.g., \texttt{(instance Pietas VirtueAttribute)}
followed by \texttt{(attribute ?AGENT Pietas)} at FOET lines
4739--4743), the auto-fix rule becomes:
\begin{quote}\small
\textbf{IF} $S \notin \mathit{BT}(X)$
\textbf{AND} $\exists$ axiom using \texttt{(instance \_ $X$)}
\textbf{AND} $\exists$ nearby \texttt{(instance $Y$ $S'$)} with
$S' \sqsubseteq S$ used as \texttt{(relation \_ $Y$)},
\textbf{THEN} propose replacing \texttt{(relation \_ $X$)} with
\texttt{(instance \_ $X$)}.
\end{quote}
This rule detected the \texttt{VirtuousAgent} instance/class
confusion in FOET and proposed the correct fix automatically.

\paragraph{AsymmetricRelation and vacuous axioms.}
Enache's GF encoding of \texttt{AsymmetricRelation} assumes
same-type arguments: \texttt{(c:Class) $\to$ (El c $\to$ El c $\to$
Formula) $\to$ Formula}.  This caused 57\% of higher-order axioms
to fail translation.  However, examining the 13
\texttt{AsymmetricRelation} instances in \texttt{Merge.kif} reveals
that 4/13 have domains on \emph{disjoint branches} (e.g.,
\texttt{attribute : Object $\times$ Attribute},
\texttt{manner : Process $\times$ Attribute}).
For cross-type relations, asymmetry is \emph{vacuously true}:
if $R : C_1 \times C_2$ with $C_1$ and $C_2$ on disjoint branches,
then $R(y,x)$ is ill-typed when $y : C_2$, $x : C_1$---the type
system prevents the symmetric case.

Thus Enache's ``failure'' to translate these axioms was
\emph{accidentally correct}: omitting a vacuously true axiom is
sound.  The proper resolution is:
\begin{itemize}[nosep]
  \item \textbf{Same-branch} (9/13): keep the axiom---asymmetry has
    real content (e.g., $\texttt{properPart}(x,y) \to
    \lnot\texttt{properPart}(y,x)$).
  \item \textbf{Disjoint-branch} (4/13): omit the axiom; the NTT verifies
    theorem that disjoint branch membership implies asymmetry.
\end{itemize}
The check is fully automatable: read \texttt{(domain R 1 C1)} and
\texttt{(domain R 2 C2)}; if $\mathit{BT}(C_1) \cap
\mathit{BT}(C_2) = \{\texttt{Entity}\}$, asymmetry is vacuous.

\paragraph{Full FOET repair results.}
Applying the NTT pipeline to the FOET seed ontology
(6\,700+ lines) detected 20 repair items:
Artifact pointers (snapshot): \texttt{Mettapedia/Languages/GF/SUMO/SumoNTT.lean},
\texttt{Mettapedia/Languages/GF/SUMO/RepairLog.lean},
\texttt{Mettapedia/Languages/GF/SUMO/data/decision\_coverage.json}.

\begin{center}
\small
\begin{tabular}{@{}rlccc@{}}
\toprule
\textbf{\#} & \textbf{Issue} & \textbf{Fixed?} & \textbf{Conf.}
  & \textbf{Auto?} \\
\midrule
 1 & Pain classification (cross-branch) & No & 0.80 & No \\
 2 & \texttt{contraryAttribute} arity & Partial & 0.95 & Yes \\
 3 & \texttt{attribute} domain (Agent$\to$Object) & Yes & 1.00 & Yes \\
 4 & AutonomousAgent (14 relations demand it) & Yes & 0.95 & Yes \\
 5 & SentientAgent hierarchy & Yes & 0.95 & Yes \\
 6 & Relation hierarchy (infra sorts) & Yes & 0.85 & Yes \\
 7--8 & Transitive closure + domain repair & Yes & 1.00 & Yes \\
 9 & AsymmetricRelation (vacuous for 4/13, verified) & Yes & 0.95 & Yes \\
10 & VirtuousAgent instance/class & Yes & 0.90 & Yes \\
11--12 & \texttt{capableInSituation} arity + domains & Yes & 1.00 & Yes \\
13 & \texttt{interferesWith} removed (fabricated) & Yes & 1.00 & Yes \\
14 & Pietas args swapped (3$\checkmark$ vs 1$\times$) & Yes & 0.98 & Yes \\
15 & Process in Object slot: \texttt{attribute}$\to$\texttt{property} & Yes & 0.92 & Yes \\
16 & \texttt{contraryAttribute} with classes$\to$schema & Yes & 0.95 & Yes \\
17 & Pain in FOET \texttt{attribute} & No & 0.85 & No \\
18 & \texttt{attribute} args swapped (6361) & Yes & 0.95 & Yes \\
19 & \texttt{holdsEthicalPhilosophy} swapped & Yes & 0.90 & Yes \\
20 & Syntax errors (14 items batch) & Yes & 1.00 & Yes \\
\bottomrule
\end{tabular}
\end{center}

In the current repair-log snapshot, 20 items are tracked; 18 are fully automatable.
The two requiring judgment
both involve Pain's cross-branch classification (Decisions 1 and 17).
The dominant repair patterns are:
\begin{enumerate}[nosep]
  \item \textbf{Counting}: correct vs incorrect usages of the same
    concept resolve ambiguity (Decisions 14, 18, 19).
  \item \textbf{Definitional unfolding}: if $X$ has a well-typed
    definitional expansion, replace the ill-typed shorthand with
    the expansion (Decision 10: VirtuousAgent $\to$
    $\exists V.\;\texttt{VirtueAttribute}(V) \land
    \texttt{attribute}(A,V)$).
  \item \textbf{Domain lookup}: read \texttt{(domain R N C)} from
    KIF and generate signatures mechanically (Decisions 3, 7, 11, 12).
  \item \textbf{Vacuity detection}: cross-type higher-order axioms
    are vacuously true by type disjointness (Decision 9).
  \item \textbf{Superrelation widening}: when $S \notin
    \mathit{BT}(X)$ for relation $R$, check if $R$ has a
    superrelation $R'$ via \texttt{(subrelation R R')} whose
    domain \emph{does} cover $X$'s sort.  Decision 15: 7 FOET axioms
    used \texttt{(attribute ?ACT ?VIRTUE)} where \texttt{?ACT} is a
    Process.  Since \texttt{(subrelation attribute property)} and
    \texttt{(domain property 1 Entity)}, the fix is
    \texttt{attribute}$\to$\texttt{property}---preserving semantics
    (every \texttt{attribute} fact entails the corresponding
    \texttt{property} fact) while widening the domain to include
    Processes.
\end{enumerate}

\paragraph{WM evidence evaluation.}
For candidates requiring semantic judgment, the repair pipeline
connects to \OSLF{} Stage~3 (world-model evidence).  For each
candidate $\varphi_i$, we construct the evidence query
$\mathit{qsem}(\lnot\varphi_i)$---``is there evidence against
$\varphi_i$?''---and evaluate it against available knowledge
(ontology axioms, encyclopedias, domain standards).  If the world
model produces a counterexample, $\varphi_i$ is too strong.  The
\emph{most general unifier} (weakest surviving candidate) wins.

\medskip\noindent\textbf{Application: Pain classification (Decision~1).}
Four candidates were evaluated:
\begin{center}\small
\begin{tabular}{@{}llcl@{}}
\toprule
& \textbf{Candidate} & \textbf{Refuted?} & \textbf{Key evidence} \\
\midrule
(a) & \texttt{PathologicProcess} (status quo) & Yes & IASP, SEP, SNOMED, SUMO's own axioms \\
(b) & \texttt{EmotionalState} (Enache) & No & But too narrow (misses sensory) \\
(c) & \texttt{StateOfMind} & No & IASP: ``psychological state''; broadest viable \\
(d) & Split: Sensation + Process & No & Clean but adds complexity \\
\bottomrule
\end{tabular}
\end{center}
Candidate (a) was refuted by 7 sources (IASP~2020 definition:
``an unpleasant sensory and emotional \emph{experience}'';
Stanford~SEP: pain as quale/mental state; SNOMED~CT:
``sensation quality''; SUMO's own documentation calling Pain
``a physical \emph{sensation}''; and SUMO's internal inconsistency
where \texttt{Medicine.kif} uses \texttt{(attribute~?P~Pain)}).
Candidate~(c) is the MGU: \texttt{StateOfMind} $\sqsubseteq$
\texttt{PsychologicalAttribute} $\sqsubseteq$
\texttt{Attribute}, making \texttt{contraryAttribute} type-correct
while being broader than \texttt{EmotionalState} to capture
pain's sensory component.

This resolves the last remaining judgment-dependent repair in the current snapshot:
of 20~decisions in the repair log, \textbf{19 are now fully
automatable}, with Pain resolved via WM evidence evaluation.
Artifact pointers (snapshot): \texttt{Mettapedia/Languages/GF/SUMO/RepairLog.lean},
\texttt{Mettapedia/Languages/GF/SUMO/DocEvidence.lean},
\texttt{Mettapedia/Languages/GF/SUMO/stratum1\_completion\_certificate.txt},
\texttt{Mettapedia/Languages/GF/SUMO/stratum1\_with\_foet\_completion\_certificate.txt}.

\paragraph{Significance.}
The full repair pipeline---NTT gap detection $\to$ candidate
enumeration $\to$ WM evidence evaluation $\to$ MGU selection---is
a closed loop connecting \OSLF{}'s three stages.  Stage~1 (GF sort
discipline) detects type gaps.  Stage~2 (behavioral types via
$\DiamondOp/\BoxOp$) computes reachability.  Stage~3 (world-model
evidence) evaluates semantic candidates.  The six repair patterns
(counting, definitional unfolding, domain lookup, vacuity detection,
superrelation widening, WM evidence evaluation) together resolve
95\% (19/20) of the detected issues fully automatically in the current repair-log
snapshot, with the
remaining 5\% (Pain) resolved semi-automatically via structured
evidence gathering from domain standards (IASP), encyclopedias
(Wikipedia, SEP), and clinical ontologies (SNOMED~CT).

\paragraph{Formal inconsistency detection via the proven-sound checker.}
The full pipeline---\SUMO{} KIF $\to$ GF Pattern $\to$ GSLT $\to$
\OSLF{} $\to$ WM \texttt{checkLang}---formally detects inconsistencies
using \SUMO{}'s own data against its own type constraints.  For each
axiom, we encode the arguments as \texttt{Pattern} trees positioned
at their KIF-declared sort (e.g., \texttt{lift\_Process(pain)} for Pain,
\texttt{lift\_Attribute(pleasure)} for Pleasure), then evaluate the
formula $\mathit{is\_S} \lor \DiamondOp(\mathit{is\_S})$ (``is on
branch $S$ or can reach branch $S$ via coercion'') using
\texttt{checkLang} with the hierarchy GSLT as the \texttt{LanguageDef}.

For \texttt{(contraryAttribute Pleasure Pain)}, which requires both
arguments to reach \texttt{Attribute}:
\begin{center}\small
\begin{tabular}{@{}llll@{}}
\toprule
\textbf{Argument} & \textbf{Starting sort} & \textbf{Reaches Attribute?}
  & \textbf{Result} \\
\midrule
Pleasure & \texttt{lift\_Attribute} & \texttt{.sat}
  & Correct \\
Pain & \texttt{lift\_Process} & \texttt{.unsat}
  & \textbf{Type gap} \\
Pain (repaired) & \texttt{lift\_Attribute} & \texttt{.sat}
  & Fixed \\
\bottomrule
\end{tabular}
\end{center}
The checker derives: $\text{arg}_1$ \texttt{.sat} $\land$ $\text{arg}_2$
\texttt{.unsat} $=$ axiom is \textbf{ill-typed}---an inconsistency
in \SUMO{}'s own type system, detected formally and proven sound via
\texttt{checkLangUsing\_sat\_sound}.  After reclassifying Pain to
\texttt{EmotionalState} (Attribute branch), both arguments return
\texttt{.sat} and the axiom becomes well-typed.

No external evidence is needed for \emph{detection}---\SUMO{}'s own
KIF data, run through the pipeline's proven-sound checker, produces
the inconsistency verdict.  External evidence (IASP, similarity-class
counting) is needed only for choosing \emph{which repair} to apply.

A comprehensive audit ran \texttt{checkLang} against all 26~repair
decisions: 20/26 are testable as reachability queries (the remaining
6 are encoding choices, meta-decisions, or syntax fixes).
All 20 testable decisions passed---confirming that each detected gap
is a genuine type error and each repair resolves it.  The audit
itself caught a test setup error (stratum~2 concept tested against a
Layer~1 GSLT), demonstrating the value of systematic verification.
The expanded audit (D22--D26) adds stratum-based checks for
LinguisticExpression, Artifact, CognitiveAgent, and Group,
with bidirectional testing (pre-fix \texttt{.unsat} $+$
post-fix \texttt{.sat}) for typed-argument-swap and
domain-narrowing archetypes.  Results are persisted as
structured WM witnesses in the Evidence Model
(\S\ref{sec:evidence-model}).
Artifact pointers (snapshot): \texttt{Mettapedia/Languages/GF/SUMO/SumoNTT.lean},
\texttt{Mettapedia/Languages/GF/SUMO/data/wm\_witnesses.jsonl},
\texttt{Mettapedia/Languages/GF/SUMO/data/decision\_coverage.json}.

\paragraph{Related work.}
This approach complements existing ontology debugging work:
\'Alvez et al.~(2019) developed white-box test generation for
FOL ontologies (applied to Adimen-\SUMO{}); Kalyanpur et al.~(2005)
introduced MUPS-based root cause analysis.  Our contribution is
the connection to \OSLF{}'s behavioral type theory: the class
hierarchy as GSLT provides a \emph{semantic} foundation for gap
detection (not just syntactic test generation), and the WM evidence
calculus provides a principled framework for candidate evaluation
(not just user interaction).  The proven-sound checker provides
a formal guarantee that ``\texttt{.sat} $\Rightarrow$ \texttt{sem}
holds'' (Theorem~\texttt{checkLangUsing\_sat\_sound}) in this pipeline
slice; in the cited ontology-debugging tooling above, we are not aware of
an equivalent theorem-level checker soundness result.

\end{document}
